\documentclass[11pt,a4paper]{article}
\usepackage[T1]{fontenc}
\usepackage[utf8]{inputenc}
\usepackage{lmodern,microtype}
\usepackage[a4paper,margin=28mm]{geometry}
\usepackage{amsmath,amssymb,amsthm,mathtools,booktabs,longtable,array,needspace}
\usepackage[dvipsnames]{xcolor}
\usepackage[colorlinks=true,linkcolor=MidnightBlue,citecolor=MidnightBlue,urlcolor=MidnightBlue]{hyperref}
\newtheorem{theorem}{Theorem}[section]
\newtheorem{proposition}[theorem]{Proposition}
\newtheorem{lemma}[theorem]{Lemma}
\newtheorem{corollary}[theorem]{Corollary}
\theoremstyle{definition}\newtheorem{definition}[theorem]{Definition}
\theoremstyle{remark}\newtheorem{remark}[theorem]{Remark}
\newtheorem{question}[theorem]{Question}
\newtheorem{conjecture}[theorem]{Conjecture}
\DeclareMathOperator{\FPdim}{FPdim}
\DeclareMathOperator{\Irr}{Irr}
\DeclareMathOperator{\Tr}{Tr}
\DeclareMathOperator{\Rep}{Rep}
\let\Vec\relax
\DeclareMathOperator{\Vec}{Vec}
\newcommand{\C}{\mathcal C}
\newcommand{\Z}{\mathcal Z}
\newcommand{\Q}{\mathbb Q}
\newcommand{\N}{\mathbb Z_{\geq0}}
\newcommand{\R}{\mathbb R}
\newcommand{\CC}{\mathbb C}
\newcommand{\ZZ}{\mathbb Z}
\newcommand{\id}{\operatorname{id}}
\allowdisplaybreaks[2]
\setlength{\emergencystretch}{2em}

\DeclareMathOperator{\Gal}{Gal}
\DeclareMathOperator{\Norm}{Norm}
\DeclareMathOperator{\disc}{disc}
\numberwithin{equation}{section}
\setcounter{tocdepth}{1}
\hypersetup{pdftitle={Grothendieck rings of complex fusion categories: rank 4},pdfauthor={Sébastien Palcoux}}
\title{Grothendieck rings of complex fusion categories:\\ rank 4}
\author{Sébastien Palcoux}
\date{}
\begin{document}
\raggedbottom
\maketitle
\begin{abstract}
Every complex fusion category of rank four admits a spherical structure.
Together with the pivotal classification of Dong, Palcoux and Plessis,
this gives exactly fifteen Grothendieck rings. We prove sphericality by
comparing the ordinary ring with the signed ring of the sphericalization.
Formal codegrees and moments of central twists bound the dimension of a
putative counterexample. An exact search leaves two signed Casimir
polynomials for one ordinary ring. Both violate the same moment inequality.
For the six rings with a non-self-dual simple, we also give a direct
classification by central induction and twist moments.
\end{abstract}
\par\smallskip
\begin{center}\begin{minipage}{0.88\linewidth}\small\itshape
GPT-6.2 Astra assisted with the research. Complete human verification is
pending, and the results should be treated with caution.
\end{minipage}\end{center}
\tableofcontents
\section{Classification and sphericality}\label{sec:introduction}
A fusion category of rank four has four simple objects up to isomorphism.
Its Grothendieck ring records the multiplicities $N_{XY}^Z$ in their tensor
products. This article is the rank-four part of a series on ranks at most four over
$\CC$. A based isomorphism preserves the unit and the basis of simple classes.

\begin{theorem}\label{thm:classification}
Every complex fusion category of rank four admits a spherical structure.
Exactly fifteen Grothendieck rings occur up to based isomorphism, as listed
in Table~\ref{finite:tab:models}. Each has a unitary realization. Moreover,
\[
 \FPdim(\C)\le24+12\sqrt3,\qquad \max_{X,Y,Z}N_{XY}^Z\le6 .
\]
Equality in either bound characterizes the near-group ring
$\operatorname{NG}(C_3,6)$.
\end{theorem}

The existence of a spherical structure concerns every category with the
given fusion rules. Exhibiting a unitary model proves existence of one
categorification for each ring. The step connecting the two assertions
is the sphericality theorem proved here (Theorem~\ref{thm:universal-spherical}).
It allows us to apply the independently established pivotal classification
\cite[Theorem~1.1 and Table~1]{DPP}.

\subsection*{History and scope}
Ostrik's rank-two and pivotal rank-three classifications established how
much information about a fusion category its center retains
\cite{OstrikRankTwo,Ostrik3}. In rank four, Rowell, Stong and Wang treated
unitary modular categories \cite{RSW}, and Bruillard classified premodular
categories up to Grothendieck equivalence \cite{Bruillard}. Hagge and Hong
observed that a pivotal category of rank at most four admits a spherical
structure \cite[Lemma~4]{HaggeHong}. Obtaining pivotality for an
arbitrary fusion category is the remaining hypothesis in that observation.

Larson reduced the pseudo-unitary case with a non-self-dual simple to seven
rings and determined the near-group parameters for $C_3$ \cite{Larson}.
Edie-Michell, Izumi and Penneys excluded the remaining quadratic candidate
under a unitarity assumption \cite{EIP}. Dong, Chen and Wang proved
sphericality for fusion categories with a subcategory of corank one and
classified the rank-four cases \cite{DCW}. Dong described the possible
self-dual tables with a proper nontrivial fusion subcategory \cite{DongFour}.
For prescribed rank and multiplicity, Vercleyen and Slingerland developed
a general enumeration of fusion rings \cite{VS}. Liu, Palcoux and Ren
classified the categorifiable multiplicity-one rings up to rank six
\cite{LPR}. Vercleyen and Slingerland subsequently computed explicit
associativity, braiding and pivotal data in low-rank multiplicity-free
examples \cite{VSData}.

Dong, Palcoux and Plessis proved the fifteen-ring theorem for pivotal
complex fusion categories, combining an a priori dimension bound with
a finite classification \cite{DPP}. Our finite calculation asks a
different question: can the signed table of a category without a
spherical structure be consistent with its ordinary table? Ruling this
out extends their theorem. The direct treatment of the non-self-dual
case in Section~\ref{sec:ns} recovers its six rings without calling on the
pivotal classification. Model constructions and their sources are collected
in Table~\ref{finite:tab:models}.

\subsection*{How the proof works}
Pass to the sphericalization $\mathcal S\to\C$ of
Etingof, Nikshych and Ostrik \cite{ENO}. Each simple object has two lifts.
Their tensor multiplicities add to give the ordinary structure constants
and subtract to give the signed ones. If signs can be chosen to make the
tables agree, the chosen lifts form a tensor subcategory equivalent to
$\C$, which inherits a spherical structure. To detect such a choice,
we use Ostrik's formal codegrees and central induction
\cite{OstrikCodegrees,Ostrik3} and the higher indicators of Ng and
Schauenburg \cite{NS,NSHigher}.

\begin{theorem}[A bound before the finite search]\label{thm:uniform-global}
Suppose a rank-four complex fusion category has no spherical structure.
Its simple objects are self-dual, its global Frobenius--Perron dimension
$D$ has degree two or four, and
\[
 D<1000,\qquad N_{XY}^Z\le31 .
\]
\end{theorem}

Three parts of the argument can be read separately.
\begin{enumerate}
\item Sections~\ref{sec:core}--\ref{sec:central-columns} establish the
comparison of ordinary and signed tables. Squared central dimensions
give a positive measure on roots of unity. Its first few moments bound
the mass contributed by induction of the unit. These estimates give
sphericality directly in several families and identify the two dimension
fields in a putative counterexample.
\item Sections~\ref{sec:quad}--\ref{sec:global-bound-proof} prove
Theorem~\ref{thm:uniform-global}. The quadratic argument compares norms
and pairs of conjugates. The cubic case gives sphericality. In degree
four a trace bound reduces the remaining problem to a finite calculation
with algebraic integers.
\item Sections~\ref{app:ordinary-census} and \ref{sec:finite-descent}
enumerate the bounded ordinary rings and test their signed lifts. The
norm of a central dimension leaves only one ordinary tuple and two signed
polynomials. The moment inequality excludes both by hand. Section~\ref{finite:sec:catalog}
then identifies the rings and gives the sharp bounds.
\end{enumerate}

Throughout, $d_i$ are the simple Frobenius--Perron dimensions and $x_i$
the signed spherical dimensions of chosen lifts. We write
$D=\sum_i d_i^2$, $E=\sum_i x_i^2$. The formal codegrees of the ordinary
and signed rings are $f_i$ and $g_i$, respectively, counted with their
character multiplicities. The two numbers that enter the final comparison are
$u=\sum_i f_i^{-2}$ and $v=\sum_i g_i^{-2}$.
\section{Sphericalization, central induction, and rigidity}\label{sec:core}
Throughout, a fusion category is finite, semisimple, rigid and
$\CC$-linear, with simple tensor unit. Its rank is the number
of isomorphism classes of simple objects. A fusion ring is a based
ring with nonnegative integral structure constants and the usual
coefficient-of-unit pairing. A \emph{Grothendieck ring} is a fusion
ring realized by such a category. The Frobenius--Perron dimension is
the positive real character $X\mapsto d_X$ of its fusion ring,
and $D=\sum_Xd_X^2$. We call a category weakly integral if
$D\in\ZZ$, and pseudo-unitary if its categorical global dimension
agrees with $D$. A pivotal structure is a monoidal isomorphism from the identity to
the double dual; it is spherical when its left and right traces agree.
We impose no pivotal hypothesis; the general existence question is
\cite[Conjecture~2.8]{ENO}.

For a self-dual basis $1,X,Y,Z$ we use the convention
\begin{equation}\label{eq:tuple-convention}
\begin{aligned}
X^2&=1+aX+bY+cZ,&XY&=bX+dY+eZ,&XZ&=cX+eY+fZ,\\
Y^2&=1+dX+gY+hZ,&YZ&=eX+hY+iZ,&Z^2&=1+fX+iY+jZ.
\end{aligned}
\end{equation}
Thus a table is denoted by $(a,b,c,d,e,f,g,h,i,j)$.
The same convention applies to signed tables. All subsequent changes
of labels or signs are explicitly local to the argument concerned.

\begin{lemma}[Rank-four algebras and coefficient bounds]
\label{core:lem:rank-four-algebra}\label{core:lem:coefficient-bound}
A semisimple complex algebra of dimension four admitting a
one-dimensional representation is commutative. In particular every
rank-four fusion ring, and every rank-four signed sphericalization
quotient, has commutative complexification. For an ordinary fusion ring,
write $d_i=\FPdim(b_i)$ and $D=\sum_jd_j^2$. Then
\[
 N_{ij}^{k}\le d_i<\sqrt D.
\]
\end{lemma}
\begin{proof}
The only noncommutative semisimple complex algebra of dimension four
is $M_2(\CC)$, which has no one-dimensional representation.
The regular based $*$-algebras in question are semisimple; their
Frobenius--Perron or signed dimension characters supply the required
representation. In the ordinary ring, reciprocity gives
$N_i^*=N_{i^*}$, so commutativity makes $N_i$ normal. Its operator norm
is its spectral radius $d_i$. Every matrix entry is bounded by that
norm, and $d_i^2<D=\sum_jd_j^2$.
\end{proof}

For a commutative self-dual table its characters have real values,
and its formal codegrees are
$f_\chi=\sum_i\chi(X_i)^2$. Orthogonality gives
\begin{equation}\label{core:orthog}
\sum_\chi\frac{\chi(X_i)\chi(X_j)}{f_\chi}=\delta_{ij},
\qquad\sum_\chi f_\chi^{-1}=1.
\end{equation}
All codegrees are positive. The $d$-number and cyclotomicity conditions
are collected in Proposition~\ref{finite:prop:standard}. In degree two,
the $d$-number criterion says that the norm divides the square of the
trace. A real character field of a categorification is totally real
and abelian.

The sphericalization $\mathcal S\to\C$ of Etingof--Nikshych--Ostrik
\cite[Remark~3.1]{ENO} is a spherical fusion category with two lifts of
each simple. Tensoring with its central sign $\varepsilon$ exchanges
these lifts. The sign has dimension $-1$ and its canonical central
lift has twist $-1$ (the sign representation in the equivariantization
description of \cite[Section~5.3]{Bartlett}). Choose lifts $X_i^+$
of dimension $x_i$ compatibly with duality. Such choices exist because
the dimensions of duals agree and the two lifts have opposite nonzero
dimensions. The group of lift signs therefore has to act simultaneously
on a dual pair. Put
\[
 R=K_0(\C)=K_0(\mathcal S)/(\varepsilon-1),\qquad
 T=K_0(\mathcal S)/(\varepsilon+1),\qquad L=\Q(x_i).
\]
The dimension character of $\mathcal S$ factors through $T$.
All $x_i$ are totally real \cite[Corollary~2.10]{ENO}.
At any chosen embedding we may change the lifts so their
embedded dimensions $\beta_i$ are positive. Write $E=\sum_i\beta_i^2$,
and use $f_i,g_i$ for the formal codegrees of $R,T$, respectively.
Set $u=\sum_i f_i^{-2}$ and $v=\sum_i g_i^{-2}$.

\begin{lemma}[Sphericalization and induction]\label{core:lem:family-induction}
For every rank-four complex fusion category, with the choices above,
\begin{enumerate}
\item its ordinary and signed coefficients satisfy
$|T_{ij}^k|\le N_{ij}^k$ and $T_{ij}^k\equiv N_{ij}^k\pmod2$;
\item all ordinary and signed codegrees $h$ belong to $L$, are
$d$-numbers, and $E/h$ is an algebraic integer;
\item $g_i\le D$ for every signed codegree. In particular $E\le D$.
If $E=D$ the positive lifts form a spherical section of $\mathcal S\to\C$.
\end{enumerate}
All statements hold at each embedding of $L$, with lift signs
chosen there. Thus a categorification admitting no spherical structure
satisfies $E<D$ at every embedding.
\end{lemma}
\begin{proof}
The two lift multiplicities have sum $N_{ij}^k$ and difference
$T_{ij}^k$. Over $\Q$, the central idempotents $(1\pm\varepsilon)/2$
split $K_0(\mathcal S)$ into $R$ and $T$. Both factors are
commutative by Lemma~\ref{core:lem:rank-four-algebra}. A character
of either factor takes values of equal absolute value on the two
lifts, so the codegrees of $\mathcal S$ are $2f_i,2g_i$.
Its spherical dimension is $2\sum_i x_i^2$. Central induction gives
the integrality and field assertions at the original spherical embedding
\cite[Theorem~2.13 and Corollary~2.15]{Ostrik3}. They persist under
every embedding since the sets of codegrees are Galois invariant.
Ostrik's $d$-number theorem applies to $2f_i,2g_i$ and multiplication
by a rational scalar preserves Galois invariance of the principal
fractional ideal.

Each signed multiplication matrix is dominated in absolute value by
the corresponding nonnegative ordinary matrix. Spectral comparison
bounds $|\chi(X_i^+)|$ by $d_i$ for every character of $T$.
Its codegree $\sum_i|\chi(X_i^+)|^2$ is thus at most $D$.
If $E=D$, then $\beta_i=d_i$ for all $i$. Subtracting the dimension
equations gives $\sum_k(N_{ij}^k-T_{ij}^k)d_k=0$. Each summand is
nonnegative, so the two tables agree. The chosen lifts form a
tensor-closed subcategory on which the forgetful functor is an
equivalence, transporting a spherical structure to $\C$.
\end{proof}

From here to the end of Section~\ref{sec:moments} the ordinary ring
is self-dual. Its chosen lifts are then also self-dual.

The central quotients $E/f_i,E/g_i$ belong to the dimension image of
the signed order, since the forgetful image is an integral signed
combination of the chosen lifts.
For a field algebra, the coefficient-of-unit functional satisfies
\begin{equation}\label{core:frobenius-trace}
\tau_T(a)=\Tr_{L/\Q}(a/E).
\end{equation}
Indeed $\Tr(a)=\sum_i\tau_T(x_i a x_i)=\tau_T(aE)$.
In particular $\Tr(x_i/E)=\delta_{i0}$.
When both ordinary and signed algebras are the same quartic field,
\begin{equation}\label{core:discriminants}
\disc(R)=N(D),\quad \disc(T)=N(E),\quad N(E/D)\in\ZZ_{>0}^2.
\end{equation}
The first two identities follow by comparing the trace pairing with
the orthonormal based pairing. The last follows because the ratio
of two order discriminants is a rational square and $E/D$ is integral.

\subsection{Spectral and trace saturation}

A self-dual based ring has a basis $b_0=1,b_1,b_2,b_3$ and a symmetric
integral multiplication tensor, with coefficient-of-unit form
$\langle b_i,b_j\rangle=\delta_{ij}$. A signed table below has the same
unit and form and allows negative structure constants.
A gauge change replaces each $b_i$ by $s_i b_i$, where $s_0=1$ and
$s_i\in\{1,-1\}$.

\begin{lemma}\label{trace:lem:absolute-rigidity}
Let $N=(N_{ijk})$ be the structure tensor of a self-dual rank-four fusion ring, and let $T=(T_{ijk})$ be a symmetric integral structure tensor defining an associative unital algebra on the same orthonormal basis $1,X,Y,Z$, with $T_{0ij}=\delta_{ij}$. If
\[
|T_{ijk}|=N_{ijk}\qquad(0\le i,j,k\le3),
\]
then there are signs $s_0=1$, $s_1,s_2,s_3\in\{1,-1\}$ such that
\[
T_{ijk}=s_i s_j s_k N_{ijk}.
\]
\end{lemma}
\begin{proof}
Associativity and symmetry give, for every $i,j$,
\[
\sum_k T_{iik}T_{jjk}=\sum_k T_{ijk}^{\,2}.
\]
The same identity holds for $N$, and the right-hand sides agree. Each term
$N_{iik}N_{jjk}-T_{iik}T_{jjk}$ is nonnegative; therefore each is zero. For a fixed $k$, all nonzero coefficients $T_{iik}$ consequently have the same sign.

Use the usual ten-entry convention
\[
(a,b,c,d,e,f,g,h,i,j)
=(N_{111},N_{112},N_{113},N_{122},N_{123},
N_{133},N_{222},N_{223},N_{233},N_{333}).
\]
The coefficients with a repeated index fall into three groups
\[
(a,d,f),\qquad(b,g,i),\qquad(c,h,j).
\]
Independent sign changes of $X,Y,Z$ make every signed coefficient in these groups nonnegative. Hence the two tables now agree except possibly at $e$. If they agree there too, we are done. Otherwise $e>0$ and the signed table has coefficient $-e$.

Compare the associativity equations for the two choices of sign of $e$. Their linear and constant parts give
\begin{gather*}
a=d+f,\qquad g=b+i,\qquad j=c+h,\\
bh+ci=bc,\qquad dh=cd+fh,\qquad bf+di=fi,\\
e^2=1+df+bi+ch.
\end{gather*}
Interchanging $Y,Z$ if necessary, assume $d\ge f$. If $d=f=0$ the
first group is zero. If $d=f>0$, the two middle relations involving
$d,f$ give $b=c=0$, whence $e^2=1+d^2$, impossible. If $d>f$,
then $bf=i(f-d)$ gives $i=bf=0$. When $f>0$, the second group
vanishes. When $f=0$, the equation $dh=cd$ gives $h=c$ and
$e^2=1+c^2$, hence $c=h=j=0$. Thus some diagonal group is zero.

Flip the basis sign belonging to a zero group. It preserves all nine diagonal coefficients and reverses the sign of $e$. The two tensors now agree, proving the lemma.
\end{proof}

\begin{corollary}\label{trace:cor:trace-rigidity}
Under the same hypotheses, suppose only that $|T_{ijk}|\le N_{ijk}$. Let $C_N$ and $C_T$ be multiplication by the respective Casimir elements $\sum_i b_i^2$. If
\[
\operatorname{Tr}(C_T)=\operatorname{Tr}(C_N),
\]
then $T$ and $N$ are gauge equivalent.
\end{corollary}
\begin{proof}
In the regular representation the multiplication matrices are symmetric, and
\[
\operatorname{Tr}(C_N)=\sum_{i,j,k}N_{ijk}^2,
\qquad
\operatorname{Tr}(C_T)=\sum_{i,j,k}T_{ijk}^2.
\]
Equality, together with termwise domination, forces equality of every absolute value. Apply Lemma~\ref{trace:lem:absolute-rigidity}.
\end{proof}

\begin{corollary}\label{core:trace-spherical}
Equality of the ordinary and signed Casimir traces in a rank-four
self-dual sphericalization implies that the original category is spherical.
\end{corollary}
\begin{proof}
Corollary~\ref{trace:cor:trace-rigidity} supplies a basis sign change
for which the signed and ordinary tensors agree. The chosen lifts
are tensor closed, so the spherical structure restricts to a tensor
subcategory equivalent to the original category.
\end{proof}

\begin{lemma}[Galois transport]\label{core:lem:galois-transport}
If a spherical dimension character belongs to the Galois orbit of the
Frobenius--Perron character, the category has a pseudo-unitary Galois
conjugate with the same fusion rules. Its central identities can
therefore be evaluated with positive Frobenius--Perron dimensions.
\end{lemma}
\begin{proof}
Extend the required embedding of the dimension field to an automorphism
of $\CC$. Apply it to skeletal tensor, duality and spherical
structure constants. All their identities and the fusion rules persist,
while the dimension character becomes the Frobenius--Perron character.
This requires no algebraic field of definition for the associativity
coefficients. Equivalently, apply the automorphism to the central
identities; restriction multiplicities are fixed and twists remain
roots of unity.
\end{proof}

\begin{corollary}[Multiplicity-one sphericality]
\label{core:cor:multiplicity-one-spherical}
Every complex categorification of a self-dual rank-four fusion ring
of multiplicity one admits a spherical structure.
\end{corollary}
\begin{proof}
For its signed quotient, domination and parity give
$|T_{ij}^{k}|\le N_{ij}^{k}$ and
$T_{ij}^{k}\equiv N_{ij}^{k}\pmod2$. Since $N_{ij}^{k}\in\{0,1\}$,
$|T_{ij}^{k}|=N_{ij}^{k}$ for every coefficient.
Lemma~\ref{trace:lem:absolute-rigidity} makes the two tables gauge
equivalent. The corresponding tensor-closed lifts form a spherical
subcategory equivalent to the original category.
\end{proof}

\subsection{Structural inputs and central traces}

\begin{proposition}[Structural inputs]\label{finite:prop:standard}
\begin{enumerate}
\item In any complex categorification every formal codegree is a
$d$-number. For its monic minimal polynomial
$t^e+a_1t^{e-1}+\cdots+a_e$, this means
$a_e^j\mid a_j^e$ for $1\leq j\leq e$
\cite{OstrikCodegrees}.
\item If $S\subseteq R$ is a fusion subring, $D/\FPdim(S)$ is an
algebraic integer \cite[Theorem 7.17.6]{EGNO}.
\item Every character value of a complex-categorifiable fusion ring is
cyclotomic \cite[Theorem 8.51 and Corollary 8.53]{ENO}.
\item In a spherical categorification with dimension character $\chi$,
every $\chi(b_i)$ is nonzero, $\chi(b_i)=\chi(b_{i^*})$, and its global
dimension $\Delta=f_\chi$ satisfies $\Delta/f\in\overline{\ZZ}$
for every formal codegree $f$. Its dimension field contains all formal
codegrees \cite[Theorem 2.13 and Corollaries 2.14--2.15]{Ostrik3}.
\item \label{finite:lem:two-dimensions}A weakly integral fusion
category is pseudo-unitary \cite[Proposition~8.24]{ENO}.
If the simple Frobenius--Perron dimensions take at most two values,
then the category has a pseudo-unitary Galois conjugate
\cite[Lemma~3.1]{Schopieray}. Indeed the values are $1$ and
possibly a single $d>1$, precisely the hypothesis of that lemma.
These statements apply in every rank.
\end{enumerate}
\end{proposition}

For a spherical category $\mathcal A$, write
$F:Z(\mathcal A)\to\mathcal A$ for the forgetful functor and $I$
for its right adjoint. With $H=\dim(\mathcal A)$, the center formula
for higher indicators reads
\begin{equation}\label{core:central-indicator}
\sum_{A\in\Irr Z(\mathcal A)}[I(V):A]\dim(A)\theta_A^n=H\nu_n(V).
\end{equation}
Here the twists are roots of unity, and
$\nu_n(V)\in\Q(\zeta_n)$; see \cite[Theorem~4.1]{NS} and
\cite[Theorem~5.1]{NSHigher}. When the fusion ring is commutative, induction of the unit has one
multiplicity-one constituent per character, of dimension $H/f_\chi$
and twist one \cite[Theorems~2.5 and~2.13]{Ostrik3}.
\begin{lemma}[The central sign and its translate twists]
\label{core:lem:translate-twists}
Let $\mathcal S$ be a rank-four self-dual sphericalization, with
central sign $\varepsilon$, and let $A_\chi\subset I(1)$ be the
multiplicity-one constituent corresponding to a character $\chi$.
Its canonical central sign translate satisfies
\[
 \theta_{\widehat\varepsilon\otimes A_\chi}=\chi(\varepsilon).
\]
Thus ordinary constituents have translate twist $+1$ and signed
constituents have translate twist $-1$. The constituents of $I(1)$
and their sign translates are disjoint. Moreover
$\nu_n(\varepsilon)=0$ for odd $n$ and $-1$ for even $n$.
\end{lemma}
\begin{proof}
Central induction identifies modular characters by
\[
 \frac{S_{A_\chi,W}}{S_{A_\chi,0}}=\chi(F(W))
\]
\cite[Example~2.9 and the paragraph preceding Theorem~2.13]{Ostrik3}.
At $W=\widehat\varepsilon$ the ratio is $\chi(\varepsilon)$;
it is also $\dim(\varepsilon)$ times the monodromy scalar.
Since $\dim(\widehat\varepsilon)=\theta_{\widehat\varepsilon}=-1$
and $\theta_{A_\chi}=1$, balancing gives the asserted twist.
The constituents have dimensions $E/f_i,E/g_i>0$, whereas their
translates have the opposite dimensions, proving disjointness.
For the indicator, the invariant space is zero for odd tensor powers
of $\varepsilon$. For even powers it is one-dimensional, and the
rotation is multiplied by the inverse pivotal map $-1$.
\end{proof}

\subsection{The rank-two input}

We use the following part of Ostrik's classification as a subcategory
obstruction \cite[Main Theorem, p.~178]{OstrikRankTwo}.
\begin{proposition}[Ostrik]\label{core:prop:rank-two}
Let $\mathcal D$ be a complex fusion category with simple objects
$1,X$. Then $X^2=1+nX$ with $n=0$ or $1$.
\end{proposition}
The statement concerns arbitrary complex fusion categories, so it applies
before a pivotal structure has been established.
\section{Indicator inequalities}\label{sec:moments}
\begin{lemma}[Positive-measure inequalities]\label{core:lem:measure-tools}
Let $\mu$ be a finite positive measure on the unit circle, with mass
$M$ and real first two moments $A,B$. Then
\[
 0\le M+B\le2M,\qquad 2A^2\le M(M+B).
\]
If its first $m$ moments all equal the real number $c$, then
\[
 \mu(\{1\})\le\frac{M+mc}{m+1}.
\]
These statements also apply after subtracting any known positive
atomic submeasure.
\end{lemma}
\begin{proof}
Writing $z=e^{it}$, one has
$\int\cos^2t\,d\mu=(M+B)/2$; Cauchy--Schwarz gives the first
inequalities. The nonnegative Fej\'er polynomial
$F_m(z)=|1+z+\cdots+z^m|^2/(m+1)$ has value $m+1$ at one and
integral $M+mc$, proving the atom bound.
\end{proof}

\begin{lemma}[Second moment]\label{core:lem:family-signed-second}
Let $u,v$ be the ordinary and signed inverse-square codegree sums.
For any signed dimension character,
\[
 E(u+v-1)\le\sum_{i=0}^3|\beta_i|
 \le1+\sqrt{3(E-1)}.
\]
\end{lemma}
\begin{proof}
The eight constituents of $I(1)$ have positive dimensions $E/f$
and squared mass $H=E^2(u+v)$. Their translates by the canonical
central sign object have opposite dimensions, are disjoint from
them, and have twists $\pm1$. The center's total squared dimension
is $4E^2$. The central indicator formula
\cite[Theorem~4.1]{NS} gives
\[
 \sum_A (\dim A)^2\theta_A^2=4Es,
 \qquad s=1+\sum_{i=1}^3\nu_2(X_i^+)\beta_i.
\]
Here each $\nu_2(X_i^+)=\pm1$. Changing a lift reverses the pivotal
map. This negates its spherical dimension and its second indicator,
the eigenvalue of pivotal transposition on its one-dimensional
space of invariant bilinear forms. Removing the sixteen
known objects, the weighted sum of squared real parts of the
remaining twists is
\[
 \tfrac12\bigl((4E^2-2H)+(4Es-2H)\bigr)
 =2E^2+2Es-2H\ge0.
\]
This gives the first bound, and Cauchy's inequality gives the second.
\end{proof}
\begin{lemma}[Quartic trace and norm]\label{core:lem:family-trace-norm}
If $L$ has degree four, the signed rational algebra is $L$ itself,
and its four character codegrees are the conjugates of $E$, with
multiplicity. Moreover
\[
 \Tr_{L/\Q}(E)\le\sum_j f_j,
 \qquad \Norm_{L/\Q}(E/f)\in\ZZ_{>0}
\]
for every ordinary codegree $f$. If $v$ is the signed inverse-square
sum, then
\[
 v\ge\frac12-\frac92\,
 \frac{\Tr_{L/\Q}(E)}{\Norm_{L/\Q}(E)}.
\]
\end{lemma}
\begin{proof}
The dimension character is a surjection from a rational algebra of
dimension four onto $L$, hence an isomorphism. The Casimir's regular
trace is the sum of the codegrees. Its trace is also
$\sum_i\Tr(T_i^2)=\sum_{i,j,k}(T_{ij}^k)^2$, at most the
corresponding ordinary sum. Norm integrality is
Lemma~\ref{core:lem:family-induction}. For $q_i=1/g_i$ the reciprocal
identity gives $\sum q_i=1$, and the four-variable
Schur inequality $9e_3(q)+2\sum q_i^2\ge1$ gives the result,
since $e_3(q)=\Tr(E)/\Norm(E)$. For clarity, its homogeneous
form is $P=e_1^3-4e_1e_2+9e_3\ge0$. Subtract the smallest entry
$t$ from all four variables to obtain $y$ with one zero entry.
Expansion gives
$P(q)=P(y)+2t e_2(y)+3t^2e_1(y)+4t^3$, and $P(y)\ge0$ is the
ordinary three-variable Schur inequality.
\end{proof}

\subsection{Averaging the central indicators}

We record once the calculation used for both quartic fields and
near-group rings. Let $\mathcal A$ be spherical with totally real
Galois dimension field $K$ and global dimension $H$.
In a finite Galois extension containing $K$, the central twists
and the required roots of unity, let $\operatorname{Av}$ denote the
average over all automorphisms $\sigma$. Write $a_A=\dim A$ for a
central simple and $\delta_z$ for the unit mass at $z$. The measure
\[
 \mu_{\mathcal A}
 =\operatorname{Av}_\sigma\sum_{A\in\Irr Z(\mathcal A)}
     \sigma(a_A/H)^2\,\delta_{\sigma(\theta_A)}
\]
is a probability measure: $a_A=\dim F(A)$ belongs to $K$, and
$\sum_A a_A^2=H^2$. Enlarging the averaging field leaves it unchanged.
Multiplying \eqref{core:central-indicator} by $\dim V$, summing over
simple $V$, and using adjunction gives
$\sum_Aa_A^2\theta_A^n=H\sum_V(\dim V)\nu_n(V)$.
Because $\nu_n(V)\in\Q(\zeta_n)$, the two Galois averages factor
whenever $K\cap\Q(\zeta_n)=\Q$. Hence
\begin{equation}\label{core:eq:averaged-indicators}
 \int z^n\,d\mu_{\mathcal A}
 =\sum_{V\in\Irr\mathcal A}
   \frac{\Tr_{K/\Q}(\dim V/H)}{[K:\Q]}\,
   \operatorname{Av}_\sigma\sigma(\nu_n(V))
 \qquad (K\cap\Q(\zeta_n)=\Q).
\end{equation}

Suppose now that $L$ has degree four, so the signed rational algebra
is $L$. We apply this identity to the sphericalization
$\mathcal A=\mathcal S$, with $H=2E$.

\begin{proposition}\label{core:prop:quartic-individual-disjointness}
Suppose that, for every integer $1\le n\le m$,
\begin{equation}\label{core:eq:quartic-individual-intersection}
 L\cap\Q(\zeta_n)=\Q.
\end{equation}
Then
\begin{equation}\label{core:eq:quartic-averaged-bound}
 2u+v\le\frac34+\frac3{m+1}.
\end{equation}
\end{proposition}

\begin{proof}
Scale the probability measure to $\mu=4\mu_{\mathcal S}$.
In \eqref{core:eq:averaged-indicators}, normalized field traces are
computed by \eqref{core:frobenius-trace}. Every term for a lift of
a nonunit simple vanishes. The surviving objects $1,\varepsilon$
have dimensions $1,-1$, with $\Tr(1/E)=1$. Thus
$\int z^n\,d\mu=(1-\nu_n(\varepsilon))/2$.
Lemma~\ref{core:lem:translate-twists} yields
\begin{equation}\label{core:eq:quartic-constant-moments}
 \int z^n\,d\mu(z)=
 \begin{cases}1,&n\text{ even},\\[2pt]\tfrac12,&n\text{ odd}.
 \end{cases}
\end{equation}
Lemma~\ref{core:lem:translate-twists} gives the indicator signs,
the eight constituents of $I(1)$ of dimensions $E/f_i,E/g_i$,
and their eight disjoint sign translates.
Thus translates of the $R$ constituents have twist $+1$, while
translates of the $T$ constituents have twist $-1$.  These statements
are preserved by every $\sigma\in G$.  Summing the squared weights
yields the two estimates
\begin{equation}\label{core:eq:quartic-atom-bound}
 \mu(\{1\})\ge 2u+v,\qquad \mu(\{-1\})\ge v.
\end{equation}

By Cauchy--Schwarz, $v\ge1/4$.  Consequently
$\nu=\mu-\tfrac14\operatorname{Dirac}_{-1}$ is a positive measure.
Its mass is $15/4$, and its first $m$ moments are all $3/4$.
Apply Lemma~\ref{core:lem:measure-tools} to obtain
$(m+1)\mu(\{1\})\le(15+3m)/4$. Together with
\eqref{core:eq:quartic-atom-bound}, this proves the claim.
\end{proof}

\begin{proposition}[Squared moments]\label{core:squared-moments}
If $L\cap\Q(\zeta_{2j})=\Q$ for $1\le j\le a$, then
\[
u+v\le\frac12+\frac{3}{2(a+1)}.
\]
In particular $u+v\le7/8$ always, and $u+v\le4/5$ if
$\sqrt2\notin L$.
\end{proposition}
\begin{proof}
Push the preceding averaged measure forward by $z\mapsto z^2$.
Its total mass is four, its first $a$ moments are one, and its atom
at one is at least $2(u+v)$. Apply Lemma~\ref{core:lem:measure-tools}. The first three
intersections are trivial because $L$ is totally real. The fourth
can be nontrivial only if $\Q(\sqrt2)\subset L$.
\end{proof}
\begin{corollary}\label{core:quartic-short-moments}
If the signed dimension field $L$ is quartic, then
\[
 \sqrt2\notin L\ \Longrightarrow\ u+v\le\frac45,
 \qquad
 \sqrt5\notin L\ \Longrightarrow\ 2u+v\le\frac98.
\]
\end{corollary}
\begin{proof}
The first assertion is Proposition~\ref{core:squared-moments}.
For the second, the real subfield at order seven has degree three,
which cannot divide four. Thus the only nontrivial possible
intersection up to order seven is $\Q(\sqrt5)$ at order five. Apply
Proposition~\ref{core:prop:quartic-individual-disjointness} with $m=7$.
\end{proof}

\begin{lemma}\label{core:lem:sqrt5-periodic-moments}
Suppose the signed dimension field is $K=\Q(\sqrt2,\sqrt5)$.
With the preceding notation, one has
\[
 5u+4v\le\frac{15}{4}.
\]
\end{lemma}
\begin{proof}
Use the unsquared averaged central measure $\mu$ above.
The nontrivial quadratic subfields of $K$ have conductors $5,8,40$.
Hence $K\cap\Q(\zeta_n)=\Q$ whenever $5\nmid n$ and
$8\nmid n$. The same indicator calculation gives
\[
 \int z^n\,d\mu=\begin{cases}1/2&n\text{ odd},\\1&n\text{ even},\end{cases}
 \qquad 5\nmid n,\quad8\nmid n.
\]
These identities hold at every such order: extending the averaging
field leaves the measure unchanged.

Restrict $\mu$ to the fortieth roots of unity. Averaging the functions
$z^{40j}$ over $j$ gives their indicator function, since the measure
has finite root-of-unity support. Thus, for each residue $n$ divisible
by neither $5$ nor $8$, the restricted measure retains the displayed
$n$th Fourier coefficient. Subtract
$\nu=\tfrac34\operatorname{Dirac}_1+\tfrac14\operatorname{Dirac}_{-1}$.
The resulting forty weights have Fourier support only at multiples
of $5$ or $8$. By the discrete Fourier transform they are therefore
\[
 r_{ij}=A_i+B_j,\qquad i\in\ZZ/8,\quad j\in\ZZ/5,
\]
where a fortieth root with exponent $l$ corresponds to
$(i,j)=(l\bmod8,l\bmod5)$.
All cells other than $(0,0)$ and $(4,0)$ are nonnegative. The forced
atoms give
\[
 r_{00}\ge2u+v-\frac34,\qquad
 r_{40}\ge v-\frac14.
\]
The total of the forty weights is at most $4-1=3$.
For an additive array there is the exact identity
\[
 \sum_{i,j}r_{ij}-5r_{00}-3r_{40}
 =\frac12\sum_{j\ne0}r_{4j}
  +\frac54\sum_{\substack{i\ne0,4\\j\ne0}}r_{ij}\ge0.
\]
Consequently
$3\ge5(2u+v-3/4)+3(v-1/4)$, proving the assertion.
\end{proof}
\subsection{Comparison of quadratic codegrees}
\begin{lemma}[Quadratic comparison]\label{core:lem:family-comparison}
In a quadratic field let $D,D'$ be the ordinary Frobenius--Perron
codegree and its conjugate. Let $E,E'$ be a signed dimension
codegree and its conjugate, with $D>E\ge E'$. Set
$P=D^{-1}+D'^{-1}$, $q=E^{-1}+E'^{-1}$, $N=DD'$ and
$k=EE'/N$. Then $k$ is a positive integer,
\begin{equation}\label{core:quadratic-strict}
 kq<P+(k-1)/D,\qquad q\le2P/3.
\end{equation}
\end{lemma}
\begin{proof}
Norm divisibility gives $k\in\ZZ_{>0}$. Since
$D>E\ge\sqrt{kN}$, compare $x+kN/x$ at $E$ and $D$ to obtain
the strict inequality. If $k\ge2$, then $D/D'>k$ and
$1/D<P/(k+1)$, whence $q<2P/(k+1)\le2P/3$.
If $k=1$, $s=D/E>1$ is a totally positive unit, so its integral
trace is at least three. Using $E\ge E'$ gives
$q/P\le2/(s+s^{-1})\le2/3$.
\end{proof}
\section{Non-self-dual categories}\label{sec:ns}

\subsection{Spherical descent}
\begin{theorem}\label{ns:thm:ns-spherical}
Every complex fusion category of rank four with a non-self-dual simple
object admits a spherical structure.
\end{theorem}

Write the dual pair as $X,Z$ and the remaining nonunit simple as $Y$.
Use the order $1,X,Y,Z$ in both sphericalization quotients. Duality fixes two basis
vectors, so its trace on the complex algebra is two. It permutes
primitive idempotents and sends character values to their complex
conjugates. Thus exactly two characters are real. Both real and
nonreal pairs are Galois stable, since reality is equivalent to
$\chi(b)=\chi(b^*)$ for every basis vector. The nonreal pair
therefore takes values in an imaginary quadratic field. Write one
as $(1,\alpha,t,\overline\alpha)$, with $t\in\ZZ$, and put
$a=\alpha+\overline\alpha$ and $b=\alpha\overline\alpha$.
Orthogonality to its conjugate gives $1+t^2+\alpha^2+\bar\alpha^2=0$.
Its codegree $h$ satisfies
\begin{equation}\label{ns:eq:ns-nonreal}
 2b=1+a^2+t^2,\qquad h=2+a^2+2t^2.
\end{equation}
In particular $a,t$ have opposite parity and $h\ge3$.

For use when $h$ is small, reciprocity gives the following common
table. In the ordinary case it is Larson's parametrization
\cite[Proposition~2.1]{Larson}. The signed case has the same six
parameters in $\ZZ$.
\begin{equation}\label{ns:eq:ns-signed-table}
\begin{aligned}
 X^2&=P X+L_0Y+C Z,& XY&=Q X+JY+L_0 Z,\\
 Y^2&=1+J X+H Y+J Z,& YZ&=L_0 X+J Y+Q Z,\\
 Z^2&=C X+L_0 Y+P Z,& XZ&=1+P X+Q Y+P Z .
\end{aligned}
\end{equation}
In either ring let $\beta$ be a character with positive values on
this basis. If $h=3$, then $(a,t)=(\pm1,0)$, and orthogonality to
$\beta$ gives $1+a\beta(X)=0$. Thus $a=-1$, $\beta(X)=1$ and
$\alpha$ is a primitive cube root of unity. Substitution in the
products $X^2,XY,Y^2$ forces $C=J=1$ and $P=L_0=Q=0$.
If $h=4$, then $(a,t)=(0,\pm1)$ and orthogonality gives $t=-1$,
$\beta(Y)=1$. Now $\alpha=\pm i$, and the same three products
give $L_0=1$, $J=Q=H=0$ and $P=C$. In each case just one integral
parameter $\mu$ is left:
\begin{equation}\label{ns:eq:two-forms}
\begin{array}{c|c|c|c}
h & (\beta(X),\beta(Y),\beta(Z)) &\text{dimension relation}
  &\text{dimension codegree}\\ \hline
3 &(1,v,1)&v^2=3+\mu v&3+v^2\\
4 &(w,1,w)&w^2=1+2\mu w&2+2w^2 .
\end{array}
\end{equation}

\begin{lemma}\label{ns:lem:ns-fields}
Suppose that $\C$ admits no spherical structure. There is a real quadratic
field $K$ and imaginary quadratic fields $L,M$ such that
\[
 R\otimes\Q\cong K\times L,\qquad
 T\otimes\Q\cong K\times M.
\]
The formal codegrees are, respectively,
\[
 D,D',f,f\qquad\text{and}\qquad E,E',g,g,
\]
where $D=\operatorname{FPdim}(\C)$ is irrational,
$f,g\in\ZZ_{\geq3}$, and primes denote the nontrivial embedding
of $K$. We can choose the lifts and a real character $\beta$ of $T$
so that all its values are positive, its codegree is $E$, and
\begin{equation}\label{ns:eq:ns-basic}
 E\geq E',\qquad E<D,\qquad
 E/D\ \text{and}\ E/g\ \text{are algebraic integers}.
\end{equation}
Every displayed codegree is greater than one, and $f\leq D$.
\end{lemma}

\begin{proof}
If $D$ were rational, weak integrality would give a spherical
structure \cite[Proposition~8.24]{ENO}. Since only two characters
are real, the Frobenius--Perron field equals $K=\Q(D)$ of degree
two. Lemma~\ref{core:lem:family-induction} puts $D$ in the signed
dimension field, also of degree at most two. This accounts for the
two real quadratic factors. The nonreal pair in each table accounts
for its imaginary quadratic factor and its repeated integral
codegree by \eqref{ns:eq:ns-nonreal}. The reciprocal-codegree
identity and positivity give $h>1$ for every codegree $h$.
Spectral domination gives $f\le D$. The integrality, sign choices
and strict inequality in \eqref{ns:eq:ns-basic} follow from
Lemma~\ref{core:lem:family-induction}.
\end{proof}

\begin{lemma}\label{ns:lem:ns-small-codegree}
In the situation of Lemma~\ref{ns:lem:ns-fields}, put
\[
 n=DD',\qquad k=\frac{EE'}{DD'}.
\]
Then $k\in\{1,2\}$ and $g\in\{3,4\}$, with $k=2$ possible only
when $g=3$.
\end{lemma}

\begin{proof}
The norm of the totally positive algebraic integer $E/D$ makes $k$
a positive integer. Set
\[
 p=1-\frac2f=\frac1D+\frac1{D'},\qquad
 q=1-\frac2g=\frac1E+\frac1{E'}.
\]
The first inequality of \eqref{core:quadratic-strict} gives
\begin{equation}\label{ns:eq:ns-key}
 kq<p+\frac{k-1}{D}\leq1+\frac{k-3}{D}.
\end{equation}
Here we used $f\leq D$. If $k\geq3$, then $q\geq1/3$ and $D\geq4$
give the contradiction
\[
 \frac{k}{3}<1+\frac{k-3}{D}\leq\frac{k+1}{4}\leq\frac{k}{3}.
\]
Thus $k\leq2$. If $k=2$ and $g\geq4$, the left side of
\eqref{ns:eq:ns-key} is at least one, whereas its right side is
$1-1/D<1$.

If $k=1$, the quadratic comparison \eqref{core:quadratic-strict}
gives $q\le2p/3<2/3$, so $g<6$.
Equation~\eqref{ns:eq:ns-nonreal} and the opposite parities of $a,t$
rule out $g=5$. Thus $g\in\{3,4\}$.
\end{proof}

\begin{proof}[Proof of Theorem~\ref{ns:thm:ns-spherical}]
Suppose that $\C$ admits no spherical structure, and use the preceding
lemmas. The two possible signed tables are given by
\eqref{ns:eq:two-forms}. In either row $E\ge E'$ forces $\mu\ge0$.
In the second row $\mu=0$ would make the whole dimension character
rational. Hence $w=\mu+\sqrt{\mu^2+1}>2$ and
$E'=2+2/w^2<5/2$. Here $k=1$, so the totally positive unit
$r=D/E>1$ has integral trace at least three. This gives
$r\ge(3+\sqrt5)/2>E'$, contrary to $D'=E'/r>1$.

In the first row, $v'=-3/v$ and
\begin{equation}\label{ns:eq:ns-ng-norm}
 E'=3+9/v^2,\qquad E+E'=\mu^2+12,\qquad EE'=3(\mu^2+12).
\end{equation}
Suppose $\mu>0$. Integrality of $E/3$ and its norm give
$(\mu^2+12)/3\in\ZZ$, hence $3\mid\mu$ and $\mu\geq3$.
Moreover
\[
 K=\Q\bigl(\sqrt{\mu^2+12}\bigr),\qquad
 \mu^2+12=3\bigl(3(\mu/3)^2+4\bigr).
\]
The radicand has $3$-adic valuation one, so $K$ is neither
$\Q(\sqrt5)$ nor $\Q(\sqrt2)$.
From \eqref{ns:eq:two-forms}, $v>7/2$ and
$v^2=3+\mu v>27/2$, whence
\begin{equation}\label{ns:eq:ns-small-conjugate}
 E'<\frac{11}{3}<2+\sqrt3.
\end{equation}
For $k=1$, the unit $r=D/E>1$ cannot have trace three, since that
would generate $\Q(\sqrt5)$. Its trace is at least four,
so $r\geq2+\sqrt3>E'$, again giving $D'<1$.

For $k=2$, put $a_0=E/D$ and $s=2D/E$. The known algebraic integer
$a_0$ has norm two, so $s=2/a_0=a_0'$ is an algebraic integer of
norm two. It is totally positive and $s>2$. Its integral trace
$s+2/s$ is therefore greater than three. Trace four would generate
$\Q(\sqrt2)$, which is excluded; thus the trace is at least
five and $s>4>E'$. Since $D'=E'/s$, this too contradicts $D'>1$.

Finally suppose $\mu=0$. Then $E=E'=6$ and $n=DD'$ is either
$36$ or $18$. Denote the unsigned parameters corresponding to
$(C,H,J,L,P,Q)$ by $(c,e,k_0,\ell,p_0,q_0)$. Their parities are
\[
 c,k_0\ \text{odd},\qquad e,\ell,p_0,q_0\ \text{even}.
\]
For a nonreal character $(1,\alpha,t,\overline\alpha)$ of $R$,
the equations for $Y(X-Z)$ and $X^2-Z^2$ give
$t=q_0-\ell$ and $\alpha+\overline\alpha=p_0-c$.
Applying \eqref{ns:eq:ns-nonreal} to $R$ therefore gives
\[
 f=2+(p_0-c)^2+2(q_0-\ell)^2\equiv3\pmod8.
\]
But $\operatorname{Tr}(D)=n(1-2/f)$ is an integer, so $f\mid2n$.
Since $f$ is odd and $n\in\{18,36\}$, it divides nine; the
congruence and $f\geq3$ force $f=3$. For $n=36$ the trace of $D$
is $12$, giving $D=D'=6$, contrary to irrationality. For $n=18$
the trace is $6$, giving negative discriminant $36-72$,
contrary to reality. Every possibility has now been excluded.
\end{proof}

\begin{corollary}\label{ns:cor:pseudo-unitary-conjugate}
Every non-self-dual rank-four complex fusion category has a
pseudo-unitary Galois conjugate.
\end{corollary}
\begin{proof}
It is spherical by Theorem~\ref{ns:thm:ns-spherical}. Duality has
trace two on the based algebra and permutes primitive idempotents,
so exactly two characters are real. This condition is Galois stable.
If $D$ is rational, weak integrality gives pseudo-unitarity. Otherwise
the Frobenius--Perron orbit consists of both real characters: its
degree is two. The spherical dimension character is real and hence
belongs to this orbit. Apply Lemma~\ref{core:lem:galois-transport}.
\end{proof}

\subsection{Two families from the repeated codegree}

Larson first reduced the pseudo-unitary non-self-dual problem to two
one-parameter families and determined its seven candidates
\cite[Sections~3--5]{Larson}. We give the reduction and the exclusions
needed here. Our argument starts with the nonreal character, avoiding
an auxiliary parametrization of all non-self-dual rings.

\begin{lemma}\label{nsnew:two-families}
A pseudo-unitary rank-four fusion category with a non-self-dual simple
has one of the following Grothendieck rings:
\begin{enumerate}
\item the near-group ring $\operatorname{NG}(C_3,m)$, with basis
$1,g,g^2,\rho$ and relations
\[
 g^3=1,\qquad g\rho=\rho,\qquad
 \rho^2=1+g+g^2+m\rho\quad(m\ge0);
\]
\item the ring $S(m)$, with basis $1,g,x,y$, where $x^*=y$, and
\[
 g^2=1,\quad gx=y,\quad gy=x,\quad
 x^2=y^2=g+m(x+y),\quad xy=1+m(x+y)\quad(m\ge0).
\]
\end{enumerate}
\end{lemma}
\begin{proof}
Let $f$ be the repeated codegree of the nonreal characters in
\eqref{ns:eq:ns-nonreal}. Write the other two codegrees as $D,D_1$, with
$D=\operatorname{FPdim}(\C)$. Thus
\[
 D^{-1}+D_1^{-1}+2/f=1,\qquad D\ge f>2.
\]
Ostrik's pseudo-unitary inequality \cite[Theorem~2.21]{Ostrik3} gives
\begin{equation}\label{nsnew:first-moment}
 u:=D^{-2}+D_1^{-2}+2f^{-2}\le\frac12(1+D^{-1}).
\end{equation}
If $D$ is irrational, $D_1$ is its quadratic conjugate. Central
induction makes $D/f$ integral, so $DD_1\ge f^2$. Therefore
\[
 1-\frac4f+\frac4{f^2}\le u\le\frac12+\frac1{2f},
\]
which gives $f\le8$. Equality at $f=8$ would require $D=f$,
contrary to irrationality. If $D$ is rational, both real codegrees
are integers: they are either conjugates of one rational number or
belong to two rational characters. Hence $D_1\ge2$, and
\[
 \frac12\ge D_1^{-1}=1-2/f-D^{-1}\ge1-3/f
\]
gives $f\le6$. In either case $2<f<8$. The expression
$f=2+a^2+2t^2$, with $a,t$ of opposite parity, leaves only
\[
 f=3:\ (a,t)=(\pm1,0),\qquad
 f=4:\ (a,t)=(0,\pm1).
\]
Substitution in \eqref{ns:eq:ns-signed-table}, already carried out in
\eqref{ns:eq:two-forms}, gives the two stated rings. In the first row
$g=X$, $g^2=Z$, $\rho=Y$. In the second, $g=Y$, $x=X$, $y=Z$.
The remaining parameter is nonnegative in the ordinary ring.
\end{proof}

\subsection{The near-group parameter from eleven moments}

The following restriction is due to Larson
\cite[Theorem~4.11]{Larson}. The Galois averaging argument of
Proposition~\ref{core:prop:quartic-individual-disjointness} gives a
short proof using the same positive central measure.

\begin{proposition}\label{nsnew:near-group}
If $\operatorname{NG}(C_3,m)$ admits a pseudo-unitary categorification,
then $m\in\{0,2,3,6\}$.
\end{proposition}
\begin{proof}
Put $\delta=\operatorname{FPdim}(\rho)$ and $D=\operatorname{FPdim}(\C)$.
Then $\delta^2=3+m\delta$ and $D=6+m\delta$.
If $\delta$ is rational, it is an integer, and
$\delta(\delta-m)=3$ with $m\ge0$ forces $\delta=3$, $m=2$.
Otherwise put $K=\Q(\delta)$, with nontrivial embedding denoted by a
prime. The codegrees are $D,D',3,3$. A constituent of $I(1)$ has
dimension $D/3=2+(m/3)\delta$; its restriction has dimension
$a+b\delta$ for nonnegative integers $a,b$. Hence $m=3k$.
The radicand $m^2+12=3(3k^2+4)$ has three-adic valuation one, so
$K$ is neither $\Q(\sqrt2)$ nor $\Q(\sqrt5)$. These are the only
real quadratic subfields of cyclotomic fields of orders at most eleven:
the maximal real subfields at orders $7,9,11$ have odd degrees $3,3,5$,
and the other orders are immediate. Thus
$K\cap\Q(\zeta_n)=\Q$ for $1\le n\le11$.

Take $\mu=\mu_{\C}$ from \eqref{core:eq:averaged-indicators}.
Since $\Tr_{K/\Q}(1/D)=1/3$ and $\Tr_{K/\Q}(\delta/D)=0$, its
$n$th moment is
$(1+\overline\nu_n(g)+\overline\nu_n(g^2))/6$, where bars denote
normalized Galois averages. An invertible object's indicator is zero
unless its order divides $n$, and otherwise is an $n$th root of unity
\cite[Lemma~3.7]{Schopieray}. Its normalized Galois average is real
and at most one. Thus the moment equals $1/6$ when $3\nmid n$
and is at most $1/2$ when $3\mid n$.

The four constituents of $I(1)$ have twist one and give the forced atom
\[
 \mu(\{1\})\ge u:=\frac1{D^2}+\frac1{(D')^2}+\frac29
 =\frac13-\frac2{3(m^2+12)},
\]
where $DD'=3(m^2+12)$. Integrating
$P(z)=|1+z+\cdots+z^{11}|^2$, whose value at one is $144$, yields
\[
 144u\le\int P\,d\mu
 \le12+2\left(\frac{66}{6}+\frac13(9+6+3)\right)=46.
\]
It follows that $m^2+12\le48$, hence $m\le6$. Together with
$3\mid m$ and the rational case $m=2$, this proves the claim.
\end{proof}

\subsection{The second family and the complete list}

Larson obtained $m\le2$ \cite[Theorem~5.8]{Larson}.
Edie-Michell--Izumi--Penneys excluded $m=2$ for unitary
categorifications \cite{EIP}, and the pivotal classification of
Dong--Palcoux--Plessis excludes it in the pseudo-unitary setting
\cite{DPP}. The two vanishing indicators of a non-self-dual simple
give a direct exclusion.

\begin{proposition}\label{nsnew:second-family}
A pseudo-unitary categorification of $S(m)$ requires $m\le1$.
\end{proposition}
\begin{proof}
Suppose $m\ge2$, set $d=m+\sqrt{m^2+1}$ and
$D=2+2d^2=4+4md$. The codegrees are $D$, $2+2d^{-2}$, $4$, $4$.
Thus the three nonunit summands of $I(1)$ have dimensions
$d^2,1+md,1+md$. Their restrictions contain the unit once, no $g$,
and respectively $2m,m,m$ copies of $x,y$ in total: this follows
from $FI(1)=4\cdot1+2mx+2my$ and irrationality of $d$.
Consequently their combined twist-one mass in $I(x)\oplus I(y)$ is
\[
 W=2md^2+2m(1+md)=4m+6m^2d.
\]
Both first and second indicators of $x$ and $y$ vanish. Integrating
$|1+z+z^2|^2$ against the positive twist measure of
$I(x)\oplus I(y)$ therefore gives $9W\le3(2dD)$.
But
\[
 3W-2dD=4m+(2m^2-8)d>0\qquad(m\ge2),
\]
a contradiction.
\end{proof}

\begin{theorem}\label{nsnew:classification}
The non-self-dual Grothendieck rings of rank-four complex fusion
categories are precisely
\[
 \operatorname{NG}(C_3,0),\quad\operatorname{NG}(C_3,2),\quad
 \operatorname{NG}(C_3,3),\quad\operatorname{NG}(C_3,6),\quad
 S(0),\quad S(1).
\]
Every such category admits a spherical structure, and its
Frobenius--Perron dimension is at most $24+12\sqrt3<45$.
\end{theorem}
\begin{proof}
Corollary~\ref{ns:cor:pseudo-unitary-conjugate} gives a pseudo-unitary
Galois conjugate with the same ring. Apply
Lemma~\ref{nsnew:two-families} and the two preceding propositions.
For existence use the models and references in
Table~\ref{finite:tab:models}. Their dimension formulas are
$3+\delta^2$ with $\delta=(m+\sqrt{m^2+12})/2$ in the near-group
case, and $2+2d^2$ with $d=m+\sqrt{m^2+1}$ for $S(m)$.
They give the stated upper bound. Sphericality was already established
in Theorem~\ref{ns:thm:ns-spherical}, independently of these models.
\end{proof}
\section{Categories with a nontrivial invertible object}\label{sec:inv}
\begin{theorem}\label{inv:thm:rankfour-invertible-spherical}
Every self-dual rank-four complex fusion category with a nontrivial
invertible simple object admits a spherical structure.
\end{theorem}

Use the sphericalization notation of Section~\ref{sec:core},
choosing positive signed dimensions $\beta_i$ and a maximal conjugate
$E$. Lemma~\ref{core:lem:family-induction} shows that, in the absence of a spherical section,
\begin{equation}\label{inv:eq:proper-signed-strict}E<D.\end{equation}
If $D$ is rational, the
category is weakly integral and already spherical by
\cite[Propositions~8.23--8.24]{ENO}. We therefore restrict attention
to irrational $D$.

\begin{proposition}\label{inv:prop:cubic-rational-two-impossible}
Let $\C$ be a self-dual rank-four complex fusion category. Suppose
that both its ordinary rational Grothendieck algebra and the signed
rational algebra of its sphericalization are the product of $\Q$
with a totally real cubic field, and that the Frobenius--Perron and
spherical dimension characters belong to the respective cubic factors.
Then the ordinary rational character does not have formal codegree two.
\end{proposition}

\begin{proof}
Suppose its codegree is two. After relabeling, the rational character is
$(1,-1,0,0)$ on the basis $\mathbf1,G,X,Y$: orthogonality to the positive
Frobenius--Perron character fixes the negative sign and gives
$\operatorname{FPdim}(G)=1$. Thus $G$ is invertible. Its values on the
three cubic characters are all one, so $G^2=\mathbf1$, $GX=X$, and $GY=Y$.

Let $\mathcal S$ be the sphericalization, with central sign $\varepsilon$.
Choose positive lifts $G^+,X^+,Y^+$ and their sign translates.
The lift $G^+$ has dimension one, satisfies $(G^+)^2=\mathbf1$, and fixes
$X^+,Y^+$ on both sides: the alternative sign translates have the wrong
nonzero dimension. In the signed quotient write $G,X,Y$ for their
classes. Its rational character is again $(1,-1,0,0)$ and has codegree
two. Its multiplication has the form
\[
 \begin{aligned}
 X^2&=1+G+aX+bY,& XY&=bX+cY,\\
 Y^2&=1+G+cX+eY,
 \end{aligned}
\]
with integers $a,b,c,e$, which may be negative. The signed Casimir is
therefore
\[
 1+G^2+X^2+Y^2=4+2G+(a+c)X+(b+e)Y.
\]
Put $x=\dim X^+$, $y=\dim Y^+$, and $E=2+x^2+y^2$.
Since the dimension character is faithful on the cubic quotient by
$G-1$, the numbers $1,x,y$ are linearly independent over $\Q$.
Writing $p=a+c$ and $q=b+e$ gives
\begin{equation}\label{inv:eq:cubic-two-halfdimension}
 E=6+px+qy,\qquad E/2=3+(p/2)x+(q/2)y.
\end{equation}

Let $F:Z(\mathcal S)\to\mathcal S$ be forgetful and $I$ its right
adjoint. The ordinary and signed rational characters give two distinct
one-dimensional characters of $K_0(\mathcal S)$, with $\varepsilon$
acting by $+1$ and $-1$, respectively. Both have formal codegree four.
The central-induction theorem therefore supplies distinct constituents
$B_+,B_-$ of $I(\mathbf1)$, each occurring once and satisfying
\[
 \dim B_+=\dim B_-=\frac{\dim\mathcal S}{4}=\frac E2,
 \qquad \theta_{B_+}=\theta_{B_-}=1.
\]
Moreover $I(\mathbf1)$ and $I(\varepsilon)$ have no common constituent:
by adjunction and the induction formula their inner product is the
trace of multiplication by $\varepsilon$, which is zero since it
interchanges every pair of lifts. Hence each $F(B_\pm)$ contains the
unit once and contains no $\varepsilon$.

Write its signed multiplicity at $G$ as
\[
 r_G=[F(B_\pm):G^+]-[F(B_\pm):\varepsilon G^+].
\]
Comparing dimensions with \eqref{inv:eq:cubic-two-halfdimension}, using
the independence of $1,x,y$, gives $1+r_G=3$. Thus
$[F(B_\pm):G^+]\ge2$, so each $B_\pm$ occurs in $I(G^+)$ at least twice.

On the other hand,
\[
 \dim\operatorname{End}(I(G^+))
 =\sum_{V\in\operatorname{Irr}(\mathcal S)}
       \dim\operatorname{Hom}(VG^+,G^+V)=8.
\]
Here $G^+$ is invertible and commutes with every simple class; each
summand equals one. The two multiplicities already contribute
$2^2+2^2=8$. Consequently
\[
 I(G^+)=2B_+\oplus2B_-.
\]
Its twist is the identity, and its twist trace is $2E=\dim\mathcal S$.
The center formula of \cite[Theorem~4.1]{NS} for the first Frobenius--Schur indicator would give
$\nu_1(G^+)=1$. This contradicts $\nu_1(G^+)=0$, since $G^+$ is a
nontrivial simple object.
\end{proof}

\begin{proof}[Proof of Theorem~\ref{inv:thm:rankfour-invertible-spherical}]
Assume there is no spherical section of the sphericalization. We have
already excluded rational $D$. Since every simple is self-dual, the
group of invertible simples has exponent two. If it has four elements
the category is pointed and spherical. Thus it has exactly two
elements $1,G$, and $G$ either fixes or interchanges the other simples
$X,Y$.

First suppose $GX=Y$. The simple Frobenius--Perron dimensions
then have the two values $1,d$, with $d>1$. Schopieray's theorem
\cite[Lemma~3.1]{Schopieray} supplies a pseudo-unitary Galois
conjugate. Its canonical spherical structure transports back along
the inverse field automorphism, contradicting the assumption.

Now suppose $GX=X$ and $GY=Y$. The ordinary fixed-action
table for a self-dual category is Dong's
\cite[Proposition~3.2]{DongFour}. Its signed analogue again has the
character $(1,-1,0,0)$ on $1,G,X,Y$: the positive lift of $G$ has
dimension one, squares to the unit and fixes the positive lifts of
$X,Y$, since their negative lifts have the opposite dimensions.
Equivalently, the line spanned by $1-G$ realizes this character in
the signed algebra. Its codegree is two. The signed dimension field
therefore has degree at most three, and contains the irrational $D$.

If it is quadratic, the general signed codegree-two obstruction,
Theorem~\ref{quad:thm:quadratic-signed-codegree-two}, gives a spherical
structure. Its proof uses the signed table, the codegree and indicator
bounds of Sections~\ref{sec:core}--\ref{sec:moments}, and the quartic
order criterion of Corollary~\ref{col:cor:quartic-order-casimir-degree}.
The latter uses the ordinary dimensional grading and
Lemma~\ref{col:lem:rank-four-grading}, independently of the present
theorem. Finally, if the signed
field is cubic, it contains the irrational $D$, so both ordinary and
signed rational algebras are products of $\Q$ and a cubic
field, with the dimension characters in the cubic factors. Their
rational characters have codegree two. This contradicts
Proposition~\ref{inv:prop:cubic-rational-two-impossible}.
Every case gives a spherical section, completing the proof.
\end{proof}
\subsection{The Fibonacci normal form}\label{sec:fib}
Dong gives the ordinary table for a self-dual rank-four fusion
category containing a Fibonacci subcategory
\cite[Proposition~3.3]{DongFour}. The argument below also proves
that the Fibonacci relation itself forces self-duality of the
remaining simples and determines the signed table.
\begin{lemma}[Fibonacci normal form]\label{fib:lem:normal-form}
A rank-four fusion ring containing $f$ with $f^2=1+f$ is self-dual
and has a basis $1,f,X,Y$ with
\begin{equation}\label{fib:eq:fib-auto-rules}
\begin{gathered}
 fX=Y,\qquad fY=X+Y,\qquad X^2=1+aX+bY,\\
 XY=f+bX+(a+b)Y,\qquad
 Y^2=1+f+(a+b)X+(a+2b)Y,
\end{gathered}
\end{equation}
where $a,b\ge0$ are integers. Its positive character is
$(1,\phi,x,\phi x)$, where $\phi=(1+\sqrt5)/2$ and
$x^2=1+(a+b\phi)x$. In a normalized signed sphericalization quotient,
the same formulas hold with integers $\alpha,\beta$ satisfying
$|\alpha|\le a$, $|\beta|\le b$, $\alpha\equiv a$ and
$\beta\equiv b\pmod2$.
\end{lemma}
\begin{proof}
On the two outside simples, multiplication by $f$ has a symmetric
nonnegative integral matrix $M$ satisfying $M^2=I+M$. Its characteristic
polynomial is $t^2-t-1$ by irreducibility, so the diagonal entries are
$0,1$ and the off-diagonal entry is one. The ring is commutative by
Lemma~\ref{core:lem:rank-four-algebra}. Duality commutes with $M$ and
cannot exchange the two different diagonal entries. Hence all simples
are self-dual. Frobenius reciprocity and associativity give the displayed products.
For a signed quotient, first change the sign of $f$ so that
$f^2=1+f$. A sign change of an outside basis element normalizes its
nonzero off-diagonal action coefficient to one. The same calculation
then gives the signed formulas; domination and parity give the stated
restrictions on $\alpha,\beta$.
\end{proof}

\section{Central columns and descent of the dimension field}\label{sec:central-columns}

Let $\C$ be a self-dual rank-four complex fusion category,
$\mathcal S$ its sphericalization, and $\varepsilon$ its sign object.
Write $R=K_0(\C)$ and
$T=K_0(\mathcal S)/(\varepsilon+1)$, and put
$D=\operatorname{FPdim}(\C)$ and
$2E=\dim(\mathcal S)$. The ordinary and signed rational algebras
are commutative, as explained above. Let $L$ be the field generated
by the signed spherical dimensions. It is totally real and abelian.
Central induction gives $E/D\in L$, so $D\in L$.

We use the canonical normalized modular matrix
$S=s/(2E)$ of $Z(\mathcal S)$. For a one-dimensional character
$\chi$ of $K_0(\mathcal S)$, let $A_\chi$ be the corresponding
simple constituent of $I(\mathbf1)$. The central-induction theorem
and modular-character formula give
\begin{equation}\label{col:eq:central-column-character}
 [I(\mathbf1):A_\chi]=1,\qquad
 \dim A_\chi=\frac{2E}{f_\chi},\qquad
 \frac{S_{A_\chi,V}}{S_{A_\chi,0}}=\chi(F(V)).
\end{equation}
In this display $f_\chi$ is the codegree in $K_0(\mathcal S)$,
twice the corresponding codegree of $R$ or $T$.
Here $F:Z(\mathcal S)\to\mathcal S$ is forgetful; see
\cite[Example~2.9 and Theorem~2.13]{Ostrik3}.

\begin{proposition}[Simultaneous Galois columns]\label{col:prop:simultaneous-columns}
Let $L$ be the signed dimension field and $\sigma\in\Gal(L/\Q)$.
Let $B_\sigma$ be the induction constituent associated with the signed
dimension character followed by $\sigma$. If $\chi$ is any ordinary
or signed character and $f_\chi$ its codegree in $R$ or $T$, then
\[
\chi(F(B_\sigma))=\frac{f_\chi}{\sigma(f_\chi)}.
\]
Consequently the codegrees determine both its folded and signed
forgetful classes, even if the rational algebras are products of fields.
Their coefficient vectors $w_\sigma,z_\sigma$ satisfy
\[
w_{\sigma,0}=z_{\sigma,0}=1,\qquad
w_{\sigma,i}\ge |z_{\sigma,i}|,\qquad
w_{\sigma,i}\equiv z_{\sigma,i}\pmod2.
\]
\end{proposition}
\begin{proof}
Let $A_\chi$ be the induction constituent for $\chi$. Its dimension
is $E/f_\chi$ and $S_{A_\chi,0}=1/(2f_\chi)$. The modular-character
formula gives
\[
S_{B_\sigma,A_\chi}
=\frac{\sigma(E/f_\chi)}{2\sigma(E)}
=\frac1{2\sigma(f_\chi)}.
\]
Symmetry of $S$ gives the asserted character value. All characters
together separate the commutative semisimple algebra, so these values
determine each class. If the two lift multiplicities at $X_i$ are
$m_i^+,m_i^-$, the two coefficients are their sum and difference.
Adjunction gives unit multiplicity one. The sign multiplicity is zero:
$I(\varepsilon)=\widehat\varepsilon\otimes I(1)$ has negative-dimensional
constituents, whereas the constituents of $I(1)$ have positive dimensions.
\end{proof}

\begin{proposition}[Positive Galois columns]\label{col:prop:positive-Galois-columns}
Suppose $[\Q(D):\Q]=4$. Identify $R$ with
its Frobenius--Perron image and $T$ with its spherical-dimension
image in the common field $K=L=\Q(D)$. For every
$\sigma\in\operatorname{Gal}(K/\Q)$, the element
\[
 u_\sigma=\frac{D}{\sigma(D)}
\]
belongs to $R$ and has nonnegative integral coefficients in its
distinguished basis. Write $w_\sigma$ for its coefficient vector and
$z_\sigma$ for the signed coefficient vector of
$v_\sigma=E/\sigma(E)\in T$. Then
\[
 w_{\sigma,0}=z_{\sigma,0}=1,\qquad
 w_{\sigma,i}\ge |z_{\sigma,i}|,\qquad
 w_{\sigma,i}\equiv z_{\sigma,i}\pmod2.
\]
\end{proposition}
\begin{proof}
Apply Proposition~\ref{col:prop:simultaneous-columns} to all ordinary
and signed characters. In a quartic field the dimension maps are
injective, so the folded and signed elements are respectively
$D/\sigma(D)$ and $E/\sigma(E)$. The coefficient assertions are
those of the proposition.
\end{proof}

\begin{proposition}[Descent of the dimension field]\label{col:prop:dimension-field-grading}
If $L\ne\Q(D)$, then $\C$ admits a nontrivial
faithful grading by a finite cyclic group.
\end{proposition}
\begin{proof}
Choose a nonidentity $\sigma\in\operatorname{Gal}(L/\Q(D))$.
Define $B_\sigma$ as in the preceding proof. The same
calculation applies without any injectivity assumption on the
dimension maps and gives
\[
 \operatorname{FPdim}(F(B_\sigma))=D/\sigma(D)=1.
\]
Adjunction gives $[F(B_\sigma):\mathbf1]=1$. Consequently
$F(B_\sigma)=\mathbf1$ in $\mathcal S$. The character defining
$B_\sigma$ differs from the dimension character, whose constituent
is the tensor unit. Thus $B_\sigma\ne\mathbf1$ in $Z(\mathcal S)$.

A central structure on the tensor unit is a monoidal natural
automorphism $\gamma$ of $\operatorname{Id}_{\mathcal S}$, and here
$\gamma$ is nonidentity. Its value at $\varepsilon$ is one.
Indeed the modular-character formula evaluated at the canonical
central lift $\widehat\varepsilon$ gives $-1$; on the other hand
this value is $\dim(\varepsilon)$ times the monodromy scalar,
namely $-\gamma_\varepsilon$. Hence
$\gamma_\varepsilon=1$.

The scalar values of $\gamma$ therefore agree on the two lifts of
every simple of $\C$. They descend to a nonidentity
monoidal natural automorphism of $\operatorname{Id}_{\C}$:
if an ordinary fusion coefficient is nonzero, one of its lifted
coefficients is nonzero and supplies the required product relation.
The finitely many scalar values form a subgroup of $\CC^\times$,
since tensor products give products and duals give inverses. This
finite group is cyclic and gives the asserted faithful grading.
\end{proof}

\begin{lemma}\label{col:lem:rank-four-grading}
A rank-four complex fusion category with a nontrivial faithful
grading is either weakly integral or has a nontrivial invertible
simple object.
\end{lemma}
\begin{proof}
The neutral component has rank one, two, or three. In rank one,
$X\otimes X^*$ is a sum of units for every simple $X$, and its
unit multiplicity is one. Thus every simple is invertible.

In rank two, the neutral component is pointed or has Fibonacci
fusion rules by Proposition~\ref{core:prop:rank-two}. In the pointed case its global
dimension is two, and equality of the dimensions of homogeneous
components makes the whole category weakly integral. In the
Fibonacci case write $f^2=1+f$ and
$\phi=(1+\sqrt5)/2$. A component of rank one is impossible:
the integer multiplicity for multiplication by $f$ would satisfy
$m^2=1+m$. Thus there is just one other component, of rank two.
On its simples $X,Y$, multiplication by the self-dual object $f$
has a symmetric nonnegative integer matrix $M$ satisfying
$M^2=I+M$. Its eigenvalues are the two roots of $t^2-t-1$,
so its trace is one and determinant is minus one. Up to exchanging
$X,Y$,
\[
 M=\begin{pmatrix}0&1\\1&1\end{pmatrix},\qquad
 \operatorname{FPdim}(Y)=\phi\operatorname{FPdim}(X).
\]
Equality of component dimensions now gives
$(1+\phi^2)\operatorname{FPdim}(X)^2=1+\phi^2$.
Hence $X$ is a nontrivial invertible object.

In rank three, there is a single nonneutral simple $X$, and the
grading group has order two. Each neutral simple $Y$ satisfies
$Y\otimes X=m_YX$, so $\operatorname{FPdim}(Y)=m_Y$ is integral.
The two homogeneous components have equal Frobenius--Perron
dimension; hence the total dimension is an integer. See
\cite[Proposition~8.20]{ENO} for equality of component dimensions.
\end{proof}

\begin{corollary}\label{col:cor:quartic-order-casimir-degree}
If the rational Grothendieck algebra of a rank-four complex fusion
category is a quartic field, its Frobenius--Perron global dimension
has degree four.
\end{corollary}
\begin{proof}
Its Frobenius--Perron character identifies the algebra with a
quartic field $K=\Q(d_i)$. This character is injective, so the unit
is the only simple of dimension one. If $D$ were rational, its four
conjugates, counted with multiplicity, would all be $D$.
The reciprocal-codegree identity would give $D=4$, making every
simple invertible. Thus $D$ is irrational. If $\Q(D)\subsetneq K$,
the dimensional grading of Gannon--Schopieray
\cite[Section~1.2]{GannonSchopieray} is nontrivial since its group
has order $[K:\Q(D)]$. Lemma~\ref{col:lem:rank-four-grading}
would force weak integrality or a nontrivial invertible simple,
both of which we have just excluded.
\end{proof}

\begin{corollary}\label{col:cor:rank-four-dimension-field-descent}
Every nonspherical self-dual rank-four complex fusion category
satisfies $L=\Q(D)$.
\end{corollary}
\begin{proof}
Otherwise the preceding proposition and lemma apply. Weakly integral
fusion categories are spherical by \cite[Proposition~8.24]{ENO}.
Theorem~\ref{inv:thm:rankfour-invertible-spherical} gives a spherical
structure in the remaining case, a contradiction.
\end{proof}

\section{Quadratic codegree obstructions}\label{sec:quad}
Ostrik proved that formal codegrees are $d$-numbers \cite{OstrikCodegrees},
and Schopieray studied the norms and traces of formal codegrees,
including finiteness at fixed norm \cite{SchopierayNorm}. Here the
ordinary and signed codegrees occur in the same quadratic field.
Comparing their two norm pairs first eliminates rational signed rows
and then bounds the ordinary global dimension. The same list of short
quadratic character rows applies to both tables.

\begin{lemma}[The quadratic-unit gap]\label{quad:lem:unit-gap}
Let $0<t<1$ be a totally positive quadratic algebraic integer of norm
one. Then $t\le(3-\sqrt5)/2$. Unless its field is
$\Q(\sqrt5)$, one has $t\le2-\sqrt3$.
\end{lemma}
\begin{proof}
Its integral trace $t+t^{-1}$ is at least three. Trace three gives
$\Q(\sqrt5)$; otherwise the trace is at least four. Solve the
quadratic equation and take its smaller root.
\end{proof}

\subsection{A quadratic signed dimension field with codegree two}

\begin{theorem}\label{quad:thm:quadratic-signed-codegree-two}
Let $\C$ be a self-dual rank-four complex fusion category.
If the signed quotient of its sphericalization has a quadratic
spherical dimension field and a rational formal codegree two, then
$\C$ admits a spherical structure.
\end{theorem}

We prove that a proper signed quotient is impossible. Choose a field
embedding and lifts with positive dimension values, and let $E\ge E'$
be the two dimension codegrees of the signed quotient $T$.
If the ordinary Frobenius--Perron dimension $D$ is rational, the
category is weakly integral and hence pseudo-unitary \cite{ENO}.
Otherwise $D$ is quadratic in the signed dimension field $K$.
A proper quotient has $E<D$: equality forces equality in the spectral
comparison with the ordinary nonnegative table, giving a spherical
section. We do not assume that the ordinary Frobenius--Perron
\emph{character} field equals its codegree field; that distinction
will be addressed below.

\subsubsection{The signed table and its norm restrictions}

A character with codegree two is rational. Indeed its three nonunit
coordinates are totally real algebraic integers whose squares sum to
one at every embedding. Each nonzero coordinate has absolute norm at
least one and absolute value at most one at every embedding, so is
$\pm1$. Orthogonality to the positive row therefore puts the
codegree-two row in the form $(1,-1,0,0)$. Orthogonality makes every other character
value on $X$ equal to one. Thus multiplication by $X$ is a symmetric
integral involution. Its negative eigenspace is already spanned by
$1-X$, so it fixes $Y,Z$ and exchanges $1,X$. Consequently
\begin{equation}\label{quad:eq:quad-two-table}
 \begin{aligned}
 X^2&=1,&XY&=Y,&XZ&=Z,\\
 Y^2&=1+X+aY+bZ,&YZ&=bY+cZ,&Z^2&=1+X+cY+dZ,
 \end{aligned}
\end{equation}
with integral coefficients and $b^2+c^2=2+ac+bd$.
The other rational character is $(1,1,r,t)$, with
\[
 h=2+r^2+t^2>2,\qquad 2+ru+tv=0,
\]
where $(1,1,u,v)$ is the positive dimension character.

Put
\[
 p=D^{-1}+(D')^{-1},\quad q=E^{-1}+(E')^{-1}=\frac12-\frac1h,
 \quad k=\frac{EE'}{DD'}\in\ZZ_{>0}.
\]
Central induction gives integrality of $E/D,E/2,E/h$.
Each of the other two ordinary codegrees is at most $D$, so
$p\le1-2/D$, whether or not those characters are rational.
Comparison on the upper-root interval yields
\begin{equation}\label{quad:eq:quad-two-norm-bound}
 kq<p+\frac{k-1}{D}\le1+\frac{k-3}{D}.
\end{equation}
Moreover
\begin{equation}\label{quad:eq:quad-two-small-conjugate}
 s_0=\frac{kD}{E}=(E/D)'\quad\hbox{is integral},\qquad
 \operatorname{Norm}(s_0)=k,\qquad k<s_0<E'.
\end{equation}
Here $D'=E'/s_0>1$ because every formal codegree exceeds one.
For $k=1$, Lemma~\ref{quad:lem:unit-gap} gives $s_0\ge(3+\sqrt5)/2$;
for $k=2$, the integral trace $s_0+2/s_0$ is at least four
(since $s_0>2$), so $s_0\ge2+\sqrt2$.
If $h\ge6$, then $q\ge1/3$ and
\eqref{quad:eq:quad-two-norm-bound}, with $D\ge4$, forces $k\le2$:
for $k\ge3$ it would give
$k/3<(k+1)/4\le k/3$.

For $h\ge11$, norm divisibility gives $EE'\ge h^2$, whence
\[
 E'\le\frac{4h}{h-2+\sqrt{(h-2)^2-16}}
    <\frac{4h}{2h-5}\le\frac{44}{17}
    <\frac{3+\sqrt5}{2}.
\]
This contradicts \eqref{quad:eq:quad-two-small-conjugate}.
In particular $h\le10$. The values $h=5,8,9$ are impossible because $h-2=3,6,7$ is not
a sum of two squares. We next treat $h=4,6,7,10$ explicitly.

For $h=4$, the entries $r,t$ have absolute value one. Both positive
are incompatible with $2+ru+tv=0$. Both negative give $u+v=2$
and $E<6$, whereas $E\ge2/q=8$. Thus take $(r,t)=(1,-1)$.
Evaluation of \eqref{quad:eq:quad-two-table} gives
\[
 v=u+2,\quad(a,b,c,d)=(b-1,b,b+1,b+2),\quad
 u^2=(2b-1)u+2b+2.
\]
For $\Delta=4b^2+4b+9$, the codegrees have trace $2\Delta$ and
norm $8\Delta$. Since $\Delta$ is odd,
$\operatorname{Norm}(E/4)=\Delta/2$ is not an integer, a contradiction.

For $h=6$, take $(r,t)=(0,-2)$. Then
\[
 v=1,\quad(a,b,c,d)=(\mu,1,0,-1),\quad
 u^2=3+\mu u,\quad E=3+u^2,\quad E'=3+9/u^2.
\]
The ordering gives $\mu\ge0$. The trace and norm of $E$ are
$\mu^2+12$ and $3(\mu^2+12)$; integrality of $E/6$ forces
$6\mid\mu$. If $\mu>0$, then $E'<13/4$ and
$K=\Q(\sqrt{\mu^2+12})$ is ramified at three. Hence
$K$ is neither $\Q(\sqrt5)$ nor $\Q(\sqrt2)$.
For $k=1$ the trace of $s_0$ is at least four, so
$s_0\ge2+\sqrt3>E'$. For $k=2$ its trace is at least five, so
$s_0>4>E'$. Both are impossible.
If $\mu=0$, then $E=E'=6$ and $K=\Q(\sqrt3)$.
Norm two is impossible in $\ZZ[\sqrt3]$ modulo three,
so $k=1$. The norm-one unit $1<s_0<6$ has even trace four or six;
trace six gives the wrong field. Thus this branch leaves only
\begin{equation}\label{quad:eq:quad-two-residual-three}
 D=12+6\sqrt3.
\end{equation}

For $h=7$, integrality of $E/7$ gives $EE'=49A$ and
$\operatorname{Tr}(E/7)=5A/2$, so $A$ is even and $EE'\ge98$.
Consequently
\[
 E'\le(35-\sqrt{833})/2<13/4.
\]
This excludes $k=2$ and forces, for $k=1$, a unit of trace three,
hence $K=\Q(\sqrt5)$. The absolute values of $r,t$ are
one and two. Both negative give $u+2v=2$ and $E<6$, contrary to
$E\ge(E+E')/2\ge35/2$.
The two remaining sign choices yield, for an integer $j$,
\[
\begin{array}{c|c|c}
(r,t)&(a,b,c,d)&\text{equation for }v\\\hline
(1,-2)&(8j-5,4j-2,2j,j-1)&v^2+(1-5j)v+4j-2=0\\
(-1,2)&(8j+5,4j+2,2j,j+1)&v^2-(5j+1)v-4j-2=0.
\end{array}
\]
For $\Delta=25j^2\pm26j+9$ with the corresponding sign, the
codegree trace and norm are $5\Delta,14\Delta$.
Integrality of $E/7$ forces $7\mid\Delta$, while
\[
 25\Delta=(25j\pm13)^2+56
\]
then gives $v_7(\Delta)=1$. This field cannot be $\Q(\sqrt5)$.

For $h=10$, take $(r,t)=(2,-2)$; the other signs contradict
positivity or give $u+v=1$ and $E<3$, whereas $E\ge5$.
Evaluation gives $v=u+1$ and
$(a,b,c,d)=(b+1,b,b+2,b+1)$. With $m=b+1$,
\[
 u^2=(2m-1)u+m+1,\qquad
 \operatorname{Tr}(E)=2(4m^2+5),\quad
 \operatorname{Norm}(E)=5(4m^2+5).
\]
The norm is odd, contradicting the integrality of $E/2$.

\subsubsection{The codegrees two and three}

Only $h=3$ remains. Its rational row can be taken as $(1,1,0,-1)$,
so evaluation gives
\[
 (a,b,c,d)=(6j,2,0,1),\quad
 v=2,\quad u^2=6+6ju,\quad j\ge0,
\]
\[
 E=6+u^2,\qquad E'=6+36/u^2.
\]
The multiple of six follows from the integral traces of $E/2$ and
$E/3$.
Suppose first $j>0$. Then $E>42$, $E'<7$, and
\eqref{quad:eq:quad-two-norm-bound} gives $k\le6$.
The integral trace $t_0=s_0+k/s_0$ lies between $k+2$ and seven.
The field $K=\Q(\sqrt{3(3j^2+2)})$ is ramified at three,
so $t_0^2-4k$ has odd three-adic valuation.
The entire resulting list is
\[
 (k,t_0,K)=(1,4,\Q(\sqrt3)),\
 (1,5,\Q(\sqrt{21})),\
 (3,6,\Q(\sqrt6)),\
 (4,7,\Q(\sqrt{33})).
\]
Indeed $1\le k\le6$ and $k+2\le t_0\le7$ are finite integer
intervals; imposing odd three-adic valuation gives exactly these
four entries. The first two are impossible modulo three: they
would require $3j^2+2$ to be a square or seven times a square.
For the last, $s_0=(7+\sqrt{33})/2>25/4$, while $j\ge2$ gives
$E'<6+1/j^2\le25/4$; at $j=1$ the field is instead
$\Q(\sqrt{15})$.
The remaining possibility is
\begin{equation}\label{quad:eq:quad-two-pell}
 \begin{gathered}
 k=3,\quad K=\Q(\sqrt6),\quad
 2w^2-3j^2=2,\quad u=3j+w\sqrt6,\\
 DD'=48w^2,\qquad EE'=144w^2,\qquad
 p=\frac{w+j}{2w}.
 \end{gathered}
\end{equation}
The Pell equation makes $j$ even. Since $j=2$ would give $w^2=7$,
one has $j\ge4$, $w\ge5$, and
\[
 (j/w)^2=\frac23(1-w^{-2})\ge\frac{16}{25},\qquad p\ge\frac9{10}.
\]

Apply Lemma~\ref{core:lem:family-signed-second}, with $U_2,V_2$
the ordinary and signed inverse-square codegree sums:
\begin{equation}\label{quad:eq:quad-two-second-moment}
 E(U_2+V_2-1)\le1+\sqrt{3(E-1)}.
\end{equation}

Cauchy's inequality on the other two ordinary codegrees gives
\[
 U_2\ge p^2+\frac{(1-p)^2}{2}-\frac2{DD'}
       \ge\frac{163}{200}-\frac1{600}=\frac{61}{75}.
\]
The signed reciprocal identity gives
\[
 V_2=\frac14+\frac19+\frac1{36}-\frac2{144w^2}
       \ge\frac{233}{600}.
\]
Hence $U_2+V_2-1\ge121/600$.
But $E\ge(E+E')/2=EE'/12\ge300$, so
\eqref{quad:eq:quad-two-second-moment} bounds this by less than
$1/300+1/10=31/300<121/600$. This excludes the whole positive
Pell family without restrictions on the other ordinary characters.

If $j=0$, then $E=E'=12$, $K=\Q(\sqrt6)$.
Write $s_0=c+d\sqrt6$. The inequalities
$0<s_0'<1<s_0<12$ imply $1\le c\le6$ and $1\le d\le2$.
Then $0<c-d\sqrt6<1$ leaves exactly
$s_0=3+\sqrt6$ of norm three or $s_0=5+2\sqrt6$ of norm one.
Thus the remaining possibilities are
\begin{equation}\label{quad:eq:quad-two-residual-six}
 D=12+4\sqrt6\quad\hbox{or}\quad D=60+24\sqrt6.
\end{equation}

\subsubsection{Excluding the three ordinary dimensions}

The ordinary Frobenius--Perron character field is $K=\Q(D)$:
a quartic character field with quadratic global codegree is excluded
by Corollary~\ref{col:cor:quartic-order-casimir-degree}.

\begin{lemma}\label{quad:lem:three-ordinary-rows}
None of $12+6\sqrt3$, $12+4\sqrt6$, and $60+24\sqrt6$ can be the
codegree of a rank-four Frobenius--Perron character with values in
the respective quadratic field.
\end{lemma}
\begin{proof}
Write its nonunit coordinates as $a_i+b_i\sqrt d$ for $d=3,6$.
Orthogonality to the conjugate character gives $1+|a|^2=d|b|^2$.
Comparing coefficients in the codegree yields
\[
\begin{array}{c|ccc}
D&|a|^2&|b|^2&a\cdot b\\\hline
12+6\sqrt3&5&2&3\\
12+4\sqrt6&5&1&2\\
60+24\sqrt6&29&5&12 .
\end{array}
\]
Up to a simultaneous signed permutation, in the first case $b=(1,1,0)$.
Then $a_1+a_2=3$ and $(a_1-a_2)^2+2a_3^2=1$. In the second case
$b=(1,0,0)$, so $a_1=2$ and $a_2^2+a_3^2=1$. In the third case
$b=(2,1,0)$ and eliminating $a_2$ gives $(5a_1-24)^2+5a_3^2=1$.
In every case a nonunit coordinate vanishes, contrary to positivity.
\end{proof}

This completes the proof of
Theorem~\ref{quad:thm:quadratic-signed-codegree-two}.

\subsection{Short signed quadratic orbits}\label{sec:short-signed}
\begin{theorem}\label{qs:short-orbits}
Let $\C$ be a self-dual rank-four complex fusion category
with irrational quadratic Frobenius--Perron global dimension $D$.
Suppose its signed dimension field is $K=\Q(D)$.
If $\C$ is not spherical, its signed rational algebra is
$K\times K$. The two formal codegrees in the complementary character
orbit are genuinely irrational, and their product is at least $100$.
\end{theorem}

The proof has three steps. We enumerate short quadratic pairs once
for both tables, eliminate the rational complementary codegrees, and
then exclude the five surviving signed pairs by comparing their
dimensions with the ordinary ones.

\subsubsection{Common quadratic comparisons}

Assume throughout the proof that $\C$ is not spherical.
Choose the signed dimension orbit $E\ge E'$ and put
\[
 N=DD',\qquad H=EE'=kN,\qquad
 p=D^{-1}+D'^{-1},\qquad q=E^{-1}+E'^{-1}.
\]
Spectral domination and saturation give $E<D$, and central induction
gives $k\in\ZZ_{>0}$. Every ordinary codegree is at most $D$.
The two outside the Frobenius--Perron orbit therefore have reciprocal
sum at least $2/D$. Lemma~\ref{core:lem:family-comparison} gives
\begin{equation}\label{qs:comparison}
 p\le1-\frac2D,\qquad
 kq<p+\frac{k-1}{D}\le1+\frac{k-3}{D}.
\end{equation}
If $q\ge1/3$, then $k\le2$: for $k\ge3$, use $D\ge4$ in the
last inequality. When $k=2$ it gives $q<1/2$.
When $k=1$, the totally positive unit $r=D/E>1$ satisfies
\[
 r\ge\frac{3+\sqrt5}{2},\qquad
 \frac qp\le\frac2{r+r^{-1}}\le\frac23.
\]
In particular $q<2/3$. If $q=1/4$, then $E\ge8$ and $D>8$;
\eqref{qs:comparison} gives $k\le4$.

All simple Frobenius--Perron dimensions belong to $K$ by
Corollary~\ref{col:cor:quartic-order-casimir-degree}; this justifies
the ordinary dimension reconstructions below.

We repeatedly use the following two bounds, where $u,v$ are the ordinary
and signed reciprocal-square codegree sums:
\begin{equation}\label{qs:moments}
 u\ge p^2+\frac{(1-p)^2}{2}-\frac2N,
 \qquad u+v-1\le\frac{1+\sqrt{3(E-1)}}E.
\end{equation}
The first is Cauchy's inequality for the two complementary ordinary
reciprocals, and does not assume that their characters are quadratic.
The second is Lemma~\ref{core:lem:family-signed-second}.

\begin{lemma}[Small integral character rows]\label{quad:lem:small-rows}
For a self-dual integral Frobenius table of rank four, a rational
character cannot have codegree five, eight, or sixteen. Distinct
rational characters cannot have codegree pair $(3,4)$, $(3,12)$,
or $(9,9)$. A rational character orthogonal to a codegree-four character
has even codegree.
\end{lemma}
\begin{proof}
A rational character is an integral row $(1,a,b,c)$ with codegree
$1+a^2+b^2+c^2$. At codegree five its nonunit row is $(\pm2,0,0)$
up to permutation; evaluating the corresponding square would require
$4=1\pm2m$. For eight and sixteen, the nonunit squared norm is
$7$ modulo eight and cannot be a sum of three squares.
At codegrees $3,4,12,9$ the absolute nonunit rows are respectively
$(1,1,0),(1,1,1),(3,1,1),(2,2,0)$. The three excluded pairs cannot
have nonunit inner product $-1$. Finally a codegree-four row is odd
in every coordinate. Orthogonality modulo two makes the sum of the
other row's coordinates, hence its squared norm, even.
\end{proof}

\subsubsection{The complete short character chart}

\begin{lemma}\label{qs:chart}
Consider a pair of conjugate quadratic characters of a self-dual
integral Frobenius table of rank four, orthogonal to a character
whose four coordinates are nonzero. Suppose their formal codegree
$G$ is a $d$-number and $L=GG'<100$. Up to simultaneous signed
permutation of the three nonunit coordinates and conjugation, their
rows occur among the following possibilities. Here
\[
 \chi=(1,(t_i+s_i\sqrt d)/2)_{i=1}^3,\qquad
 \ell=G^{-1}+G'^{-1}.
\]
Signs of structure constants are unrestricted; this chart applies
equally to ordinary and signed tables. Entries marked $*$ have
rational codegree.
\[
\begin{array}{c|r@{\,}r@{\,}r|r@{\,}r@{\,}r|r|c|r}
d&\multicolumn{3}{c|}{t}&\multicolumn{3}{c|}{s}&L&\ell&\|U\|^2\\\hline
5&2&-1&-1&0&1&1&20&1/2&1\\
5&2&-1&1&0&1&1&25^*&2/5&6\\
5&-3&-1&-1&1&1&1&25&3/5&6\\
2&2&-2&-2&0&2&2&32&1/2&1\\
3&0&-4&-2&0&2&2&36&2/3&2\\
5&-3&-1&1&1&1&1&45&1/3&2\\
5&-3&-1&-4&1&1&2&45&2/3&2\\
13&2&-3&-3&0&1&1&52&1/2&1\\
2&2&-2&2&0&2&2&64^*&1/4&3\\
2&-2&-2&-6&2&2&4&64&3/4&3\\
2&-4&-2&0&2&2&2&72&1/3&2\\
5&2&-4&-4&0&2&2&80&1/2&1\\
3&2&-4&0&0&2&2&96&1/4&3\\
3&-4&-4&-6&2&2&4&96&3/4&3
\end{array}
\]
The last column gives the squared norm of the primitive integer
vector parallel to the nonunit part of the idempotent for this pair.
In the four rows where it is one, that basis element squares to one.
\end{lemma}
\begin{proof}
Put $J=|s|^2$ and $C=t\cdot s$. Orthogonality of the conjugate rows gives
\begin{equation}\label{qs:row-identities}
 |t|^2=dJ-4,\quad
 G=\frac{dJ+C\sqrt d}{2},\quad
 L=\frac d4\bigl(4J+|t\times s|^2\bigr),\quad
 \ell=\frac{dJ}{L}.
\end{equation}
Character values are algebraic integers, so $t_i,s_i$ are integers;
they are both even if $d\not\equiv1\pmod4$, and have the same
coordinatewise parity otherwise. Orthogonality to a nonzero-coordinate
character rules out support one for $s$. The reciprocal identity gives
$0<\ell<1$, hence $dJ<L<100$. The $d$-number condition is
$L\mid(dJ)^2$.

Here is a bounded arithmetic check of the chart. Take $s$ nonnegative
and ordered, and take $t_i\ge0$ wherever $s_i=0$.
If $d\not\equiv1\pmod4$, write $s=2b,t=2a$. The bounds are
\[
\begin{gathered}
d\in\{2,3,6,7,10,11\},\qquad
2\le|b|^2\le\lfloor24/d\rfloor,\\
|a|^2=d|b|^2-1,\qquad
|b|^2+|a\times b|^2\le\lfloor24/d\rfloor.
\end{gathered}
\]
If $d\equiv1\pmod4$, they are
\[
 d\in\{5,13,17,21,29,33,37,41\},\qquad
 2\le J\le\lfloor99/d\rfloor,
 \qquad |t|^2=dJ-4.
\]
For each permitted $d,J$, take the nonnegative ordered three-square
representations of $J$. For each such $s$, take
\[
 |t_1|\le\sqrt{dJ-4},\quad
 |t_2|\le\sqrt{dJ-4-t_1^2},\quad
 t_3=\pm\sqrt{dJ-4-t_1^2-t_2^2},
\]
retaining integral values with the stated parity. Apply
\eqref{qs:row-identities}, $dJ<L<100$ and $L\mid(dJ)^2$.
The resulting $L$ lists, before identifying the repeated values, are
\[
\begin{array}{c|l}
d=2&32,64,64,72\\
d=3&36,96,96\\
d=5&20,25,25,45,45,80\\
d=13&52,
\end{array}
\]
and the other displayed radicands give no row. Substitution gives
exactly the fourteen $t,s$ pairs above; the remaining identifications are precisely signed permutations
and conjugation. This is a finite sum-of-squares check,
with every range and divisibility specified independently of any
fusion-ring multiplicity.

The idempotent $P$ projecting onto the character pair has unit coefficient
$\ell$ and nonunit coefficients
\[
 P_i=\frac{d(Jt_i-Cs_i)}{2L}.
\]
Clear their denominators and divide by their common gcd to obtain $U$.
If $U=\lambda(P-\ell)$, its two eigenvalues are
$\lambda(1-\ell)$ and $-\lambda\ell$; this computes the last column.
In the norm-one rows $\ell=1/2$ and $\lambda=2$, so $U^2=1$.
\end{proof}

\subsubsection{Rational complementary signed codegrees}

\begin{lemma}\label{qs:rational-complement}
The two complementary signed codegrees cannot both be rational.
\end{lemma}
\begin{proof}
A rational signed codegree two already gives sphericality by
Theorem~\ref{quad:thm:quadratic-signed-codegree-two}; hence assume all
rational codegrees under consideration are at least three.

First suppose the complementary characters form a quadratic orbit,
with a repeated rational codegree $g$. Then $q=1-2/g$.
The preceding comparisons leave only $g=3,4,5$, and $k=1$ for
$g=4,5$. Lemma~\ref{qs:chart} applies with $L=g^2$. Its only constant-codegree
entry with $g\le5$ has $g=5$.

For $g=5$ the two rows, after signed permutation, are
\[
 (1,-1,\phi,\phi'),\qquad(1,-1,\phi',\phi),
 \qquad \phi=\frac{1+\sqrt5}{2},\quad\phi'=1-\phi.
\]
Evaluation of the symmetric table gives
\begin{equation}\label{quad:eq:constant-five-normal-form}
(t,t,t,t-1,t,t-1,t,t-1,t-1,t),\qquad t\in\ZZ.
\end{equation}
This calculation uses only the displayed character rows and the
orthonormal based pairing, so it also applies to ordinary tables.
The remaining characters are $(1,z+1,z,z)$, where
$z^2+(2-3t)z-t=0$. With $\Delta=9t^2-8t+4$ their codegrees satisfy
\[
 E+E'=3\Delta,\qquad EE'=5\Delta.
\]
The value $t=0$ gives rational selected dimensions. Otherwise
$\Delta\ge5$, and the smaller root is
\[
 E'=\frac{10}{3+\sqrt{9-20/\Delta}}
 \le\frac{15-5\sqrt5}{2}<\frac{3+\sqrt5}{2}\le r.
\]
This contradicts $D'=E'/r>1$.

It remains to treat two rational characters, of integral codegrees
$3\le g\le h$. The small-row restrictions are
Lemma~\ref{quad:lem:small-rows}. Since $q=1-1/g-1/h\ge1/3$, we have $k\le2$.
The conditions $q<1/2$ for $k=2$ and $q<2/3$ for $k=1$ now give
\[
\begin{array}{c|l}
k=2&(g,h)=(3,3),\\
k=1&(3,3),\ (3,h)\ (h\ge6),\ (4,4),\ (4,6),\ (4,10).
\end{array}
\]
For $(3,h)$ with $h\ge6$, central integrality gives $H\ge h^2$.
As $q=2/3-1/h$, the smaller root satisfies
\[
 E'\le\frac2{q+\sqrt{q^2-4/h^2}}
 \le9-3\sqrt5<\frac{3+\sqrt5}{2}\le r.
\]
The second bound is attained at $h=6$: both $q$ and
$q^2-4/h^2=4/9-4/(3h)-3/h^2$ increase thereafter.
For $(4,6)$, integrality of $E/4,E/6$ implies that of $E/12$ by
Bezout, so $E\ge12$ and $E'\le2$. For $(4,10)$ the same argument
gives $E\ge20$ and $E'\le5/3$. Again $D'=E'/r<1$.

For $(4,4)$ the rational rows are, up to signed permutation,
$(1,-1,1,-1)$ and $(1,-1,-1,1)$. Their orthogonal complement has
selected row $(1,1,z,z)$. Character evaluation gives
\[
 X^2=1,\quad XY=Z,\quad XZ=Y,\quad
 Y^2=Z^2=1+t(Y+Z),\quad YZ=X+t(Y+Z).
\]
Here $z^2=1+2tz$. Select the larger codegree and change both signs
if necessary, so $t\ge1,z>1$; $t=0$ would give rational dimensions.
Then $E'=2+2/z^2<5/2<r$.

Finally, for $(3,3)$ the rows are
$(1,0,1,-1)$ and $(1,1,-1,0)$, up to signed permutation.
Evaluation gives the tuple
\[
(m-1,m-1,m,m,m,m,m,m,m+1,m+1),\qquad m\in\ZZ,
\]
and remaining characters $(1,z-1,z,z+1)$ with
\begin{equation}\label{quad:eq:rational-three-normal-form}
 z^2=3mz+1,\qquad E=3+3z^2,\qquad EE'=9(9m^2+4).
\end{equation}
The value $m=0$ gives a zero simple dimension. Otherwise the larger
codegree has $|z|>3$ and $E'=3+3/z^2<10/3$.
The field $K=\Q(\sqrt{9m^2+4})$ is neither
$\Q(\sqrt2)$ nor $\Q(\sqrt5)$, since either equality
would require a square congruent to two modulo three.
If $k=1$, the positive unit $r$ cannot have trace three, so
$r\ge2+\sqrt3>E'$. If $k=2$, put $s=2D/E>2$.
This is a totally positive integer of norm two. Its trace is at least
four, and equality would give $K=\Q(\sqrt2)$. Thus its trace
is at least five, so $s>4>E'$. In either case $D'<1$, a contradiction.
\end{proof}

The selected character gives a quadratic quotient $K$ of the semisimple
rational algebra $T_{\Q}$, leaving a complementary algebra of
dimension two. If the latter is $\Q\times\Q$, both codegrees
are rational, contrary to Lemma~\ref{qs:rational-complement}. Otherwise
it is a quadratic field $K'$. Its codegree $G$ belongs to $K'$ by its
definition and to $K$ by central induction. If $K'\ne K$, then
$G\in K\cap K'=\Q$, giving a repeated rational codegree,
again excluded by the lemma. Thus $G$ is genuinely irrational and
$K'=K$. This proves the $K\times K$ assertion in
Theorem~\ref{qs:short-orbits} without asserting field containment
for arbitrary complementary character values.

\subsubsection{A signed basis involution}

\begin{lemma}\label{qs:signed-involution}
Under the common quadratic hypotheses of Theorem~\ref{qs:short-orbits},
if the signed quotient has a nonunit basis element $X$ with $X^2=1$,
then the original category is spherical.
\end{lemma}
\begin{proof}
Multiplication by $X$ is a symmetric integral involution, hence a signed
permutation matrix. It interchanges $1,X$. On the other two basis
elements it either fixes them or interchanges them. Since every selected
simple dimension is nonzero, in the fixed case the two signs equal the
selected value of $X$. Changing its basis sign gives $XY=Y,XZ=Z$;
then $(1,-1,0,0)$ is a codegree-two character, and
Theorem~\ref{quad:thm:quadratic-signed-codegree-two} applies.

In the interchange case choose signs to obtain
\[
 X^2=1,\quad XY=Z,\quad XZ=Y,\qquad
 Y^2=Z^2=1+aY+bZ,\quad YZ=X+bY+aZ.
\]
The selected dimensions can be taken as $(1,1,z,z)$, with
$z^2=1+(a+b)z$. If $z$ is rational, it is $\pm1$; central induction
then makes the ordinary global dimension rational, and weak integrality
gives sphericality. Otherwise select the larger signed codegree and
change both signs so $z>1$. Thus $z$ is a norm-minus-one quadratic
unit and $E=2+2z^2$. Here $q=1/2$, so \eqref{qs:comparison} forces
$k=1$. Properness would give a totally positive unit $r=D/E>1$.

Let $\eta>1$ be the fundamental unit of $K$. The existence of $z$
of norm minus one implies $N(\eta)=-1$. Write
$\eta-\eta^{-1}=t\in\ZZ_{>0}$. If $t\ge2$, then $z\ge\eta$, and
$r\ge\eta^2\ge3+2\sqrt2>3>E'=2+2/z^2$, which is impossible.
If $t=1$, then $K=\Q(\sqrt5)$ and $\eta=\phi$.
Write $z=\phi^{2j+1}$, $j\ge0$. For $j\ge1$,
$E'\le2+2\phi^{-6}<\phi^2\le r$. For $j=0$, only
$r=\phi^2$ keeps $D'=E'/r>1$; it gives $D=10+4\sqrt5$ and
$D^{-1}+D'^{-1}=1$. The other ordinary reciprocal codegrees would
then sum to zero. Thus no proper signed dimension exists.
\end{proof}

Apply the chart to the complementary signed orbit. The starred rows
are excluded by Lemma~\ref{qs:rational-complement}, the norm-one rows
by Lemma~\ref{qs:signed-involution}, and every row with
$q=1-\ell\ge2/3$ by \eqref{qs:comparison}. Only the following cases remain:
\begin{equation}\label{qs:five-cases}
\begin{array}{c|ccccc}
L&25&36&45&64&96\\\hline
K&\Q(\sqrt5)&\Q(\sqrt3)&\Q(\sqrt5)&
\Q(\sqrt2)&\Q(\sqrt3)\\
q&2/5&1/3&1/3&1/4&1/4.
\end{array}
\end{equation}

\subsubsection{Ordinary dimensions needed at the small endpoints}

\begin{lemma}\label{qs:ordinary-endpoints}
None of the following can be the ordinary global dimension of the
nonspherical category under consideration:
\[
 5\phi^2,\quad 12+6\sqrt3,\quad15+6\sqrt5,\quad
 75+33\sqrt5,\quad8+4\sqrt2,\quad24+16\sqrt2,\quad36+20\sqrt3.
\]
\end{lemma}
\begin{proof}
The value $12+6\sqrt3$ is already excluded by Lemma~\ref{quad:lem:three-ordinary-rows}.
For the other values the simple dimensions belong to their indicated quadratic field.
For $d=2,3$, write the three nonunit dimensions as $a_i+b_i\sqrt d$.
Orthogonality of this character to its conjugate gives
$1+|a|^2=d|b|^2$. If $D=A+B\sqrt d$, then
\[
 |b|^2=\frac A{2d},\qquad |a|^2=\frac A2-1,
 \qquad a\cdot b=\frac B2.
\]
For $d=5$, use $(t_i+s_i\sqrt5)/2$ and
\eqref{qs:row-identities}. The resulting small integer-vector equations
are recorded below. In the middle column the three entries are
$(|s|^2,|t|^2,t\cdot s)$ for $d=5$ and
$(|b|^2,|a|^2,a\cdot b)$ otherwise.
\[
\begin{array}{c|c|l}
D&\text{invariants}&\text{dimension row or consequence}\\\hline
5\phi^2&(3,11,5)&(1,\phi,\phi,\phi^2)\\
15+6\sqrt5&(6,26,12)&(1,\phi,\phi^2,\phi^3)\\
75+33\sqrt5&(30,146,66)&(1,\phi^2,\phi^3,\phi^5)\\
8+4\sqrt2&(2,3,2)&(1,1,\eta,\eta)\\
24+16\sqrt2&(6,11,8)&(1,\eta,\eta,\eta^2)\\
36+20\sqrt3&(6,17,10)&(1,\epsilon,\epsilon,2\epsilon-1).
\end{array}
\]
Here $\eta=1+\sqrt2$ and $\epsilon=2+\sqrt3$.
We justify uniqueness in this small table. For $d=2,3$, a vector
of squared norm two is a signed permutation of $(1,1,0)$, and one
of squared norm six is a signed permutation of $(1,1,2)$.
The indicated dot products and companion squared norms then force
respectively
\[
 a=(1,1,\pm1),\quad(1,1,3),\quad(2,2,3),
\]
up to the corresponding signed permutations. For $d=5$, squared
norms three, six and thirty force $s=(1,1,1),(1,1,2),(1,2,5)$.
Parity and the dot products force $t=(1,1,3),(1,3,4),(3,4,11)$.
For the last case, the particularly short check is
$|t-(11/5)s|^2=4/5$. Restoring the signs required by positive
Frobenius--Perron dimensions gives the displayed rows.

An object of dimension one is invertible, so the $8+4\sqrt2$ row gives
sphericality by Theorem~\ref{inv:thm:rankfour-invertible-spherical}.
In either of the $5\phi^2$ and $15+6\sqrt5$ rows,
an object $X$ of dimension $\phi$ satisfies $X^2=1+Y$ for some object
$Y$ of dimension $\phi$. If $Y\ne X$, reciprocity puts $X$ in $XY$;
the remaining dimension is one but $[XY:1]=0$, an impossibility.
Thus $Y=X$. Apply the normal form in
Lemma~\ref{fib:lem:normal-form}. For $5\phi^2$ its remaining
parameter is $x=\phi$, so $a+b\phi=1$, giving $(a,b)=(1,0)$;
Corollary~\ref{core:cor:multiplicity-one-spherical} gives sphericality.
For $15+6\sqrt5$ one has $x=\phi^2$, so
$a+b\phi=\phi^2-\phi^{-2}=\sqrt5=-1+2\phi$.
This forces $a=-1$, impossible already for a fusion ring.

For $75+33\sqrt5$, label the dimensions of $X,Y,Z$ by
$\phi^2,\phi^3,\phi^5$. Evaluation of
$X^2=1+aX+bY+cZ$ gives
$a+b+3c=1$ and $a+2b+5c=3$, hence $a+c=-1$.
This contradicts nonnegativity.

For $24+16\sqrt2$, $d_X=d_Y=\eta$ and $d_Z=\eta^2=1+2\eta$.
Evaluation of $X^2,Y^2$ gives $c=h=0$, $a+b=d+g=2$.
Evaluation of $XY=bX+dY+eZ$ then gives $e=1$, $b=d=0$.
Consequently $X^2=1+2X$, excluded by
Proposition~\ref{core:prop:rank-two}.
For the last row, $\epsilon^2=4\epsilon-1$ and $d_Z=2\epsilon-1$.
Evaluation of the two squares forces $X^2=Y^2=1+2Z$, so
$b=d=0$. Evaluation of $XY=eZ$ would simultaneously require
$e=1$ and $2e=4$, a contradiction.
\end{proof}

\subsubsection{Excluding the five signed rows}

For a totally positive quadratic $d$-number in
$\Q(\sqrt2)$, $\Q(\sqrt3)$ or $\Q(\sqrt5)$,
its norm has square class in, respectively,
\begin{equation}\label{qs:norm-classes}
 \{1,2\},\qquad\{1,6\},\qquad\{1,5\}.
\end{equation}
Indeed its conjugate quotient is a totally positive unit. The least
such units are $3+2\sqrt2$, $2+\sqrt3$ and $\phi^2$. Separate the
even and odd powers of that quotient: the original number is a rational
multiple of a unit or of one plus that unit. The norms of the latter
three elements have square classes two, six and five. This proves
\eqref{qs:norm-classes}, including the restrictions it imposes on the
integral ratio $k=H/N$.

The quadratic unit descriptions used below are elementary Pell
reductions. In $\ZZ[\sqrt3]$ the prime over three is generated by
$\sqrt3$ and the prime over two by $1+\sqrt3$; in
$\ZZ[\sqrt2]$ the prime over two is generated by $\sqrt2$.
Divide a generator by powers of the least positive unit until it lies
in the fundamental interval. The norm equation in that interval leaves
precisely the generators just stated. In $\Q(\sqrt5)$ the units
are $\pm\phi^n$. These assertions can also be checked directly by
writing an algebraic integer and its conjugate as $z,N(z)/z$.

\paragraph{Complementary norm $25$.}
The row in the chart gives, by its two orthogonality equations,
selected dimensions $(1,1,z,-1-z)$. Orthogonality to the selected
conjugate gives $\operatorname{Tr}(z)+2N(z)=-3$. Put
$h=N(z)+1\in\ZZ$ and $A=4h^2+5$. Direct substitution gives
\[
 E=A+2|h|\sqrt A,\qquad H=5A.
\]
The selected dimension field is $\Q(\sqrt5)$ precisely when
$A=5w^2$. This implies $5\mid h$. If $h\ne0$, then $|h|\ge10$,
since $h=5v$ gives $w^2=20v^2+1$ and $v=\pm1$ is impossible.
Thus $E\ge405$. Here $q=2/5$ gives $k\le2$, and
\eqref{qs:norm-classes} excludes $k=2$. Hence $r=D/E\ge\phi^2$.
Since $D'=E'/r>1$, we have $E'>\phi^2$, so
\[
 E<\frac1{2/5-\phi^{-2}}<56,
\]
a contradiction. If $h=0$, then $E=5$ and $r=\phi^2$ is the only
positive unit with $E'/r>1$. It gives $D=5\phi^2$, excluded by
Lemma~\ref{qs:ordinary-endpoints}.

\paragraph{Complementary norm $36$.}
The short row gives $(1,x,1,-1)$ with $N(x)=-3$.
With $\epsilon=2+\sqrt3$, select the larger square to write
\[
 x=\sqrt3\epsilon^j,\qquad E=3+3\epsilon^{2j},\qquad j\ge0.
\]
As $q=1/3$, the norm classes force $k=1$. Then $r=D/E\ge\epsilon$.
For $j\ge1$, $E'\le3+3\epsilon^{-2}<\epsilon$, impossible.
For $j=0$, $E=6$, and only $r=\epsilon$ keeps $D'>1$.
It gives $D=12+6\sqrt3$, excluded by Lemma~\ref{qs:ordinary-endpoints}.

\paragraph{Complementary norm $45$.}
The selected dimensions are $(1,1-z,-1-z,z)$, with $N(z)=-1$.
Thus
\[
 |z|=\phi^{2j+1},\qquad E=3+3\phi^{4j+2},\qquad j\ge0.
\]
Again $q=1/3$ and the norm classes give $k=1$.
Since $E'\le3+3\phi^{-2}<\phi^4$, properness forces $r=\phi^2$.
It follows that
\[
 p=\frac{\phi^2}{3}-\frac{\sqrt5}{E}.
\]
The cases $j=0,1$ give respectively $D=15+6\sqrt5$ and
$D=75+33\sqrt5$, already excluded. For $j\ge2$,
$E>300$, $N\ge1125$, and $p\ge5/6$. Here
$v=23/45-2/N$, so \eqref{qs:moments} gives
\[
 u+v-1\ge\frac{79}{360}-\frac4N>\frac15,
 \qquad
 u+v-1<\frac1{300}+\sqrt{\frac3{300}}<\frac15,
\]
a contradiction.

\paragraph{Complementary norm $64$.}
Orthogonality gives, after signs, the selected row
$(1,1-z,1+z,1)$ with $N(z)=-2$. Put $\epsilon=3+2\sqrt2$.
Then
\[
 z=\sqrt2\epsilon^j,\quad E=4+4\epsilon^{2j},\quad
 H=64w^2,\qquad j\ge0,
\]
where $w$ is the rational coefficient of $\epsilon^j$.
Here $q=1/4$, so $k\le4$; the norm classes leave $k=1,2,4$.
For $k=1$, $r\ge\epsilon$. If $j\ge1$ then $E'<\epsilon$.
If $j=0$, $E=8$ and only $r=\epsilon$ is possible; this gives
$D=24+16\sqrt2$, already excluded.

For $k=4$, the totally positive integer $s=4D/E>4$ has norm four.
Its ideal is $(2)$, so $s=2\epsilon^a$ with $a\ge1$.
But $s\ge2\epsilon>8\ge E'$, contradicting $D'=E'/s>1$.
For $k=2$, the corresponding integer $s=2D/E>2$ has norm two,
so $s=(2+\sqrt2)\epsilon^a$. Since $E'\le8$, only $a=0$ is possible.
Consequently
\[
 p=\frac{2+\sqrt2}{4}-\frac{2\sqrt2}{E}.
\]
At $j=0$, this gives $D=8+4\sqrt2$, already excluded.
For $j\ge1$, $E>100$, $H\ge576$, $p\ge5/6$ and $N=H/2$.
Since $v=19/32-2/H$, \eqref{qs:moments} gives
\[
 u+v-1\ge\frac{29}{96}-\frac6H\ge\frac7{24},
 \qquad
 u+v-1<\frac1{100}+\sqrt{\frac3{100}}<\frac{37}{200},
\]
again impossible.

\paragraph{Complementary norm $96$.}
The selected row is $(1,1-z,1+z,-1)$, so $E=4+2z^2$
and $N(z)=-2$.
Put $\epsilon=2+\sqrt3$. Selecting the larger square gives
\[
 z=(1+\sqrt3)\epsilon^j,\quad E=4+4\epsilon^{2j+1},\quad
 H=96w^2,\qquad j\ge0,
\]
where $w$ is the $\sqrt3$ coefficient of $z$.
The norm classes and $k\le4$ leave $k=1,4$.
For $k=4$, the integer $s=4D/E>4$ has norm four and equals
$2\epsilon^a$, $a\ge1$. This exceeds
$E'\le12-4\sqrt3$, a contradiction. Thus $k=1$.
The same bound for $E'$ is less than $\epsilon^2$, so $r=\epsilon$,
and
\[
 p=\frac\epsilon4-\frac{2\sqrt3}{E}.
\]
At $j=0$, $D=36+20\sqrt3$, already excluded.
For $j\ge1$, $E>200$, $N\ge864$, $p\ge11/12$, and
$v=29/48-2/N$. Hence
\[
 u+v-1\ge\frac{43}{96}-\frac4N>\frac25,
 \qquad
 u+v-1<\frac1{200}+\sqrt{\frac3{200}}<\frac25.
\]
This excludes the last row of \eqref{qs:five-cases} and proves
Theorem~\ref{qs:short-orbits}.

\subsection{The quadratic dimension bound}\label{quad:opt:section}

Suppose throughout that the category is nonspherical and $D$ is irrational
quadratic. By Corollary~\ref{col:cor:rank-four-dimension-field-descent}
the signed dimension field is $\Q(D)$. Keep $p,q,N$ from
\eqref{qs:comparison}, write $EE'=hN$, put $t=E/D$, and let $M,L$
be the ordinary and signed complementary codegree products.
Theorem~\ref{qs:short-orbits} gives $L\ge100$. With
$F(z)=z^2+(1-z)^2$ and $u,v$ as in \eqref{qs:moments},
\[
 N=\frac D{p-1/D},\qquad
 q=\frac th p+\frac{1-t^2/h}{E},\qquad E^2\ge hN.
\]
These identities allow the same second-moment comparison for both
types of ordinary character orbit.

\begin{lemma}[All norm ratios greater than one]\label{quad:opt:norm-ratio}
If $h\ge2$, then $D<1000$, without any assumption on the ordinary
complementary codegrees.
\end{lemma}
\begin{proof}
Suppose $D\ge1000$. Since $0<t<1$, its integral trace satisfies
$t+h/t\ge h+2$. Consequently
\[
 \frac th\le\frac2{h+2+\sqrt{h^2+4}}
 \le1-\frac1{\sqrt2}<\frac3{10}.
\]
We have $u\ge p^2+(1-p)^2/2-2/N$ by \eqref{qs:moments} and
$v=F(q)-2/(hN)-2/L$. Also $E>\sqrt{2D}>44$. If $E<60$, then
$p>2D/E^2>5/9$, and
\[
 q\le\frac E{2D}+\frac1E\le\frac7{150}.
\]
For the last inequality, the function $E/2000+1/E$ increases on
$[\sqrt{2000},60]$ and has terminal value $7/150$.
Thus
\[
 u+v-1\ge\frac{11}{27}+F(7/150)-1-\frac{23}{1000}
 >\frac{29}{100}
 >\frac1{44}+\sqrt{\frac3{44}}
 >\frac{1+\sqrt{3(E-1)}}E,
\]
a contradiction.

Assume now $E\ge60$. Set $a=\min(E/(2D),3/10)$ and
$G(p)=p^2+(1-p)^2/2$.
Since $t/h\le a$, the formula for $q$ gives $q\le ap+1/E$.
The function $F$ decreases up to $1/2$ and has $|F'|\le2$ on
$[0,1]$, so $F(q)\ge F(ap)-2/E$. Completing the square in
$G(p)+F(ap)-1$ and using $N>D$ gives
\[
 C(a)-\frac1{50}-\frac3D
 \le\frac3E+\sqrt{\frac3E},\qquad
 C(a)=\frac{1-2a}{3+4a^2}.
\]
The function $C$ decreases on $[0,3/10]$. The following elementary
comparisons contradict this inequality:
\[
\begin{array}{c|c|c}
 E\text{ range}&\text{left side exceeds}&\text{right side is less than}\\\hline
 {[60,100]}&90/301-23/1000>11/40&11/40\\
 {[100,200]}&5/19-23/1000>21/100&21/100\\
 {[200,400]}&15/79-23/1000>7/50&7/50\\
 {[400,\infty)}&5/42-23/1000>19/200&19/200.
\end{array}
\]
In each row the bound on $a$ follows from $D\ge1000$; the last uses
$a\le3/10$. The right column follows by evaluating at the left
endpoint. This completes the proof.
\end{proof}

\begin{lemma}[Unit-norm comparison]\label{quad:opt:unit-comparison}
Suppose $h=1$ and write $U=u+2/N$. Then
\[
 U+F(tp)-1-\frac1{50}
 \le \frac{3+\sqrt{3(E-1)}}{E}.
\]
If $E/E'\ge\kappa$ and $D\ge D_0$, then
$E\ge\max\{D_0t,\sqrt{D_0\kappa/p}\}$.
\end{lemma}
\begin{proof}
Here $q=tp+(1-t^2)/E$, and convexity gives
$F(q)\ge F(tp)+(4tp-2)(1-t^2)/E$.
In $u+v=U+F(q)-4/N-2/L$, use
$4E/N=4t(p-1/D)<4tp$. The combined error, multiplied by $E$, is at most
\[
 4tp+(2-4tp)(1-t^2)=2-2t^2(1-2tp)\le2.
\]
Indeed $0<t\le(3-\sqrt5)/2<1/2$ and $0<p<1$.
Now apply \eqref{qs:moments}. The last bound uses
$E=tD$ and $E^2=(E/E')N$ with $N>D/p$.
\end{proof}

\begin{theorem}\label{quad:opt:unconditional}
Every nonspherical self-dual rank-four complex fusion category
with irrational quadratic Frobenius--Perron global dimension satisfies
$D<1000$. Every fusion coefficient is at most $31$.
\end{theorem}
\begin{proof}
Suppose $D\ge1000$. An invertible nonunit would give sphericality
by Theorem~\ref{inv:thm:rankfour-invertible-spherical}.
Lemma~\ref{quad:opt:norm-ratio} leaves $h=1$.
We describe the possible lower bounds $H(t)$ for the left side of
Lemma~\ref{quad:opt:unit-comparison} and then check all of them together.

\emph{Large complementary product.}
If $M\ge100$, then $U=F(p)-2/M$. Minimizing
$F(p)+F(tp)-1$ in $p$ by completing the square gives
\[
 H(t)=\frac{(1-t)^2}{2(1+t^2)}-\frac1{25}.
\]
In the bound for $E$ we may replace $p$ by $1$.

\emph{Short quadratic orbit.}
If $M<100$ and the two ordinary complementary characters are conjugate,
Lemma~\ref{qs:chart} applies. The $M=20$ row forces an invertible.
For the other rows, except for the repeated codegree five, substitution
into $U=F(p)-2/M$ gives the following lower bounds $U_0$.
\[
\begin{array}{c|rrrrrr}
p&1/4&1/3&2/5&1/2&2/3&3/4\\\hline
U_0&19/32&1/2&11/25&7/16&23/45&19/32
\end{array}
\]
For the repeated codegree five, $(p,U_0)=(3/5,11/25)$.
In each case use $H(t)=U_0+F(pt)-1-1/50$.

\emph{Two rational characters.}
Write their codegrees $g\le j$, with $gj<100$. A codegree-two rational
row forces an invertible by orthogonality to the positive dimension
row. Thus $g\ge3$. Discard $5,8,16$, the pairs $(3,4),(3,12),(9,9)$,
and odd $j$ when $g=4$, by Lemma~\ref{quad:lem:small-rows}.
The pair $(4,4)$ also forces an invertible: of two distinct orthogonal
sign rows, their sum has unit coordinate two and exactly one nonzero
nonunit coordinate, of absolute value two. Orthogonality to the
positive row forces the corresponding dimension to be one.
Call the remaining set $\mathcal P$; it contains $52$ pairs. For each, put
\[
 p=1-1/g-1/j,\qquad U_0=p^2+1/g^2+1/j^2,
 \qquad H(t)=U_0+F(pt)-1-1/50 .
\]
We may take $\kappa=1$ except for $(3,3)$, where
\eqref{quad:eq:rational-three-normal-form} gives
$N=9(9m^2+4)$ with $m\ne0$. This lies strictly between consecutive
squares after division by nine. Hence $EE'=N$ is not a square,
and $E/E'$ is a nontrivial totally positive quadratic unit by
Ostrik's $d$-number theorem \cite{OstrikCodegrees}. We may take
$\kappa=5/2$ by Lemma~\ref{quad:lem:unit-gap}.

There are two ranges for the unit $t$: either $t\le2-\sqrt3<268/1000$,
or $t=(3-\sqrt5)/2<382/1000$. The latter occurs in $\Q(\sqrt5)$.
For rational complementary characters the discriminant of the ordinary
order is $Ngj$ and has square class five. By \eqref{qs:norm-classes},
$gj$ is a square or five times a square. Among $\mathcal P$ this
leaves $(3,3)$ and
\[
 \mathcal P_5=\{(3,15),(3,27),(4,20),(6,6),(7,7)\}.
\]
The field in the $(3,3)$ case is $\Q(\sqrt{9m^2+4})$, which cannot be
$\Q(\sqrt5)$ by reduction modulo three.

\emph{Endpoint check.}
All the functions $H$ just listed decrease on the indicated intervals.
The right side of Lemma~\ref{quad:opt:unit-comparison} decreases for
$E\ge2$. The table gives a subdivision for each family. Its entries are
the successive right endpoints in thousandths, starting at zero.
For $[a,b]$ in this subdivision set
\[
 Z=\max\left\{D_0a,\,
  \frac{\lfloor100\sqrt{D_0\kappa/p}\rfloor}{100}\right\}.
\]
In every interval the following two rational comparisons hold:
\[
 ZH(b)-3>0,\qquad (ZH(b)-3)^2>3(Z-1).
\]
Since $E\ge Z$, these comparisons contradict the lemma.
\begin{center}\small
\begin{tabular}{@{}lrl@{}}
\toprule
Input & $D_0$ & Right endpoints multiplied by $1000$\\\midrule
$M\ge100$ & 1000 & $60,190,360,382$\\
Six short quadratic rows & 1000 & $90,220,340,382$\\
Repeated five, $t\le2-\sqrt3$ & 1000 & $60,130,230,268$\\
Repeated five, all $t$ & 3000 & $150,340,382$\\
$\mathcal P$, $t\le2-\sqrt3$ & 1000 & $46,52,75,140,220,268$\\
$\mathcal P_5$, all $t$ & 1000 & $100,240,340,370,382$\\
\bottomrule
\end{tabular}
\end{center}
For instance, for $(g,j)=(3,6)$ on $[0,46/1000]$, the input is
$p=1/2$, $U_0=7/18$, $Z=4472/100$. Substitution gives
$ZH(b)-3>287/25$ and $3(Z-1)<263/2$.
The table is reproducible by two nested loops over its stated rows
and endpoints. The accompanying exact checker uses rational arithmetic
and integer square roots.

Only the repeated five with $t=(3-\sqrt5)/2$ and $1000\le D<3000$
remains. In \eqref{quad:eq:constant-five-normal-form}, denote the
integer parameter by $m$. The choice of basis signs does not affect
the codegrees. We have $N=5\Delta$, $D+D'=3\Delta$ and
$\Delta=9m^2-8m+4$. The $d$-number condition $N\mid(3\Delta)^2$
gives $m\equiv1\pmod5$. Since $1/D'=3/5-1/D>1/2$,
we have $D'<2$ and $\Delta<1001$. Thus $m=-9,-4,1,6$, with
$\Delta=805,180,5,280$, respectively. The selected field
$\Q(\sqrt\Delta)$ equals $\Q(\sqrt5)$ only for $\Delta=180,5$,
but $D<3\Delta<1000$ in both cases. This proves the dimension bound.
Finally every symmetric ordinary multiplication matrix has operator
norm its Perron root $d_i\le\sqrt D<32$, so all entries are at most $31$.
\end{proof}

\section{Sphericality in the cubic case}\label{sec:cub}

The even part of $\mathrm{SU}(2)_7$ is the familiar rank-four example
with a cubic dimension field \cite[Section~5.3.10]{RSW}.
We prove that a category with a cubic dimension field already admits
a spherical structure. This allows the pivotal classification to be
applied after descent, without assuming a pivotal structure here.

The Frobenius--Perron dimension field means the field generated by
all simple Frobenius--Perron dimensions. Rank-four fusion rings are commutative by
Lemma~\ref{core:lem:rank-four-algebra}. Suppose a complex fusion categorification exists
and its dimension field $K$ has degree three. The ordinary rational algebra is then $\Q\times K$,
and character cyclotomicity makes $K$ cyclic and totally real
\cite[Corollary~8.53]{ENO}. Every character is consequently real.
The based involution acts by complex conjugation on characters, so it
is the identity, so every simple is self-dual.
Its Frobenius--Perron codegree $D$ is cubic as well, by the order
discriminant argument in Lemma~\ref{cub:lem:cubic-square-norm-schur}
below (applied just to the ordinary algebra).

Central induction places $D$ in the field generated by the signed
spherical dimensions. That field has degree at most four, so it has
degree three and equals $K$. Thus the signed rational algebra is also
$\Q\times K$, with its dimension character in the cubic factor.
We write $f$ for the codegree of the ordinary rational character.

\subsection{The trace identity and a positive measure}
Put $T_{\Q}=(K_0(\widetilde{\C})/(\varepsilon+1))\otimes\Q$ for the
signed algebra of the sphericalization. Choose dual-compatible lifts
$X_i^+$ of the four simples, with $X_0^+=\mathbf1$, and write
$X_i^-=\varepsilon\otimes X_i^+$,
$x_i=\dim X_i^+$, and $E=\sum_i x_i^2$.  Let $\rho$ denote the
rational character of $T_{\Q}$ and let
$g=\sum_i\rho(X_i^+)^2$ be its codegree.  The rational algebra of the sphericalization is the commutative direct sum $R_{\Q}\oplus T_{\Q}$, where $R=K_0(\C)$.  The Frobenius trace identity is
\begin{equation}\label{cub:eq:cubic-corrected-frobenius}
 \operatorname{Tr}_{K/\Q}(x_i/E)
       =\delta_{i0}-\frac{\rho(X_i^+)}g.
\end{equation}
To see this identity, let $\tau$ be the coefficient of the unit.
In the orthonormal signed basis, the regular trace of multiplication by
$a$ equals $\tau(a C_T)$, where $C_T$ is the Casimir.  The rational and
cubic components of $C_T$ are $g$ and $E$.  Evaluating this formula on
$X_i^+ C_T^{-1}$ gives the displayed identity.

Let $f_1,\ldots,f_4$ be the formal codegrees of $K_0(\C)$,
and let $g,E_1,E_2,E_3$ be those of the signed algebra.  Put
\[
 u=\sum_{i=1}^4f_i^{-2},\qquad
 v=g^{-2}+\sum_{i=1}^3 E_i^{-2}.
\]
As usual, these numbers are rational.  The dimension-field theorem
\cite[Corollary~2.15]{Ostrik3} places all the displayed codegrees in $K$.
Throughout, codegrees are counted with their character multiplicities.
Both reciprocal-codegree sums equal one and $f,g\ge2$.

\begin{proposition}\label{cub:prop:cubic-corrected-average}
There is a positive, finitely supported measure $\xi$ on the unit circle,
of mass three, such that
\[
 \int z^n\,d\xi(z)=
 \begin{cases}1/2,&n\text{ odd},\\1,&n\text{ even},\end{cases}
 \qquad K\cap\Q(\zeta_n)=\Q.
\]
Its atoms at $1$ and $-1$ satisfy
\begin{align}
 \xi(\{1\})&\ge\frac34(2u+v)+\frac{f-2}{f^2}
                          +\frac{g-2}{2g^2},\label{cub:eq:positive-atom}\\
 \xi(\{-1\})&\ge\frac34v+\frac{g-2}{2g^2}\ge\frac14.
 \label{cub:eq:negative-atom}
\end{align}
\end{proposition}
\begin{proof}
Let $F:Z(\widetilde{\C})\to\widetilde{\C}$ be forgetful
and $I$ its right adjoint. Use the canonical normalized center matrix
$S=s/(2E)$. The rational character $\rho$ corresponds to a self-dual
constituent $B\ne1$ of $I(1)$ by
\cite[Example~2.9 and Theorem~2.13]{Ostrik3}. Put
\[
 a_A=S_{0A}=\frac{\dim A}{2E}\in K,\qquad
 b_A=S_{BA}=\frac{\rho(F(A))}{2g}\in\Q.
\]
Orthogonality gives
\begin{equation}\label{cub:eq:cubic-corrected-row-norms}
 \sum_Aa_A^2=\sum_Ab_A^2=1,\qquad \sum_Aa_Ab_A=0.
\end{equation}
Take a cyclotomic field $\Omega$ containing the modular data. For each
indicator index used, enlarge it to contain $\zeta_n$. Write
$G=\operatorname{Gal}(\Omega/\Q)$. Its signed permutations
$\widehat\sigma$ satisfy
\begin{equation}\label{cub:eq:cubic-corrected-galois}
 \sigma(a_A)=\epsilon_\sigma(A)a_{\widehat\sigma A},\quad
 b_A=\epsilon_\sigma(A)b_{\widehat\sigma A},\quad
 \theta_{\widehat\sigma A}=\sigma^2(\theta_A).
\end{equation}
The last identity has no scalar anomaly for the canonical center
normalization \cite[Lemma~4.2 and Section~4.1]{DLN}. Define
\[
 \xi=\frac1{|G|}\sum_{\sigma\in G}\sum_A
 \bigl(3\sigma^2(a_A)^2+\sigma^2(a_A)b_A\bigr)
 \operatorname{Dirac}_{\sigma^2(\theta_A)}.
\]
Here $\sigma^2$ means composition. Equation
\eqref{cub:eq:cubic-corrected-row-norms} gives mass three.

For an allowed $n$, the restrictions of $G$ to $K$ and
$\Q(\zeta_n)$ are independent. Squaring permutes
$\operatorname{Gal}(K/\Q)\cong C_3$; no corresponding
surjectivity is required on the indicator field.
Apply \eqref{core:central-indicator} and
\eqref{cub:eq:cubic-corrected-frobenius}. The moment of the first
term of $\xi$ is
\[
 \frac{1-\nu_n(\varepsilon)}2
 -\frac1{2g}\operatorname{Av}_{\sigma\in G}
 \sigma^2\!\left(\sum_i\rho(X_i^+)
          [\nu_n(X_i^+)-\nu_n(X_i^-)]\right).
\]
The sum in parentheses is $2g\sum_Aa_Ab_A\theta_A^n$ by the same
center identity. The second term of $\xi$ cancels this residual.
The sign object has indicator $-1$ for even $n$ and zero for odd $n$:
its even invariant space is one-dimensional and rotation includes its
inverse pivotal map $-1$, whereas the odd invariant space vanishes.
This proves the moment formula.

Reindexing each summand by $W=\widehat\sigma A$ fixes the twist
$\theta_W$ and gives its individual weight
\begin{equation}\label{cub:eq:cubic-corrected-weight}
 w_W=\operatorname{Tr}_{K/\Q}(a_W^2)
              +\frac{b_W}{3}\operatorname{Tr}_{K/\Q}(a_W).
\end{equation}
In particular the measure does not depend on enlarging $\Omega$.
Let $h_W=[F(W):1]-[F(W):\varepsilon]$. The Frobenius trace identity
says $\operatorname{Tr}(a_W)=h_W/2-b_W$.
Outside $I(1)\cup I(\varepsilon)$ one has $h_W=0$, so
$w_W=\operatorname{Tr}(a_W^2)-b_W^2/3\ge0$.
For $W\subset I(1)$ put $t_W=\operatorname{Tr}(a_W)>0$.
All conjugate codegrees are positive, and $b_W=1/2-t_W$; hence
\[
 w_W\ge\frac{t_W^2}{3}+\frac{(1/2-t_W)t_W}{3}
       =\frac{t_W}{6}>0.
\]
Its sign translate has the same weight. Thus positivity holds
pointwise, without grouping atoms.

For the four ordinary constituents the $b$-values are
$(f-3)/(2f),1/(2f),1/(2f),1/(2f)$; the signed values replace $f$
by $g$. This follows from $a_A=1/(2f_\chi)$ and
$b_A=1/2-\operatorname{Tr}(a_A)$. Summing their corrections in
\eqref{cub:eq:cubic-corrected-weight} gives $(f-2)/(2f^2)$ and
$(g-2)/(2g^2)$, respectively. By
Lemma~\ref{core:lem:translate-twists}, the
ordinary sign translates have twist $+1$ and the signed translates
have twist $-1$. Summing their weights gives
\eqref{cub:eq:positive-atom}--\eqref{cub:eq:negative-atom}; the final
bound uses $v\ge g^{-2}+(1-g^{-1})^2/3$.
\end{proof}

\subsection{Periodic moments}
\begin{proposition}\label{cub:prop:cubic-periodic}
Suppose the ordinary and signed rational algebras are $\Q\times K$
as above. If $F$ is the conductor of $K$, then
\begin{equation}\label{cub:eq:cubic-periodic}
 \mathcal L:=\frac32(u+v)+\frac{f-2}{f^2}+\frac{g-2}{g^2}
          \le1+\frac2F\le\frac97.
\end{equation}
\end{proposition}
\begin{proof}
A nontrivial cyclic cubic conductor is odd and at least seven.
Since $[K:\Q]=3$, the field $K$ is disjoint from
$\Q(\zeta_n)$ whenever $F\nmid n$.
Restrict $\xi$ to the $2F$th roots and call the restriction $\xi_0$.
For $F\nmid r$, average its moments with exponents $2Fj+r$ over
$0\le j<M$ and let $M\to\infty$. The geometric average is zero off
the $2F$th roots and one on them, so $\xi_0$ has the same prescribed
$r$th moment. Its mass is at most three.
Subtract $\tfrac34\operatorname{Dirac}_1+
\tfrac14\operatorname{Dirac}_{-1}$. The residual Fourier transform
on this cyclic group is supported on $\{0,F\}$, so its weights are
constant on each parity class; call them $a,b$.
Each class has roots other than $\pm1$, at which these are weights
of the positive measure $\xi_0$. Hence $a,b\ge0$ and
\[
 1+F(a+b)\le3,\qquad
 \xi(\{1\})+\xi(\{-1\})=1+a+b\le1+\frac2F.
\]
Now use \eqref{cub:eq:positive-atom}--\eqref{cub:eq:negative-atom}.
\end{proof}

\subsection{The rational codegrees and cubic norms}
The two rational codegrees $f,g$ are integers at least two. The
ordinary value $f=2$ is excluded by
Proposition~\ref{inv:prop:cubic-rational-two-impossible}.

\begin{lemma}\label{cub:lem:cubic-square-norm-schur}
For either algebra, let $r$ be its rational codegree and $H$ a codegree
from its cubic character orbit. There are positive integers $A,s$ such
that
\[
 \operatorname{Tr}(H)=rA,\qquad
 \operatorname{Norm}(H)=r^3s^2,
\]
and $H/r$ has minimal polynomial
\begin{equation}\label{cub:eq:cubic-normalized-polynomial}
 z^3-Az^2+(r-1)s^2z-s^2.
\end{equation}
If $r\ge3$, the inverse-square sum of the four codegrees satisfies
\begin{equation}\label{cub:eq:cubic-uniform-schur}
 r^{-2}+\sum_{\sigma:K\hookrightarrow\R}\sigma(H)^{-2}
 \ge r^{-2}+\tfrac12(1-r^{-1})^2.
\end{equation}
\end{lemma}

\begin{proof}
First $H$ has degree three. If it were rational, the reciprocal-codegree
identity would give $(r-1)H=3r$, so $(r,H)=(2,6)$ or $(4,4)$.
The order discriminant is the product of its four formal codegrees.
The first case gives $432$, which is not a square, contrary to the
cyclic cubic field discriminant. The second gives $256$; its square
root would have to be divisible by the conductor $F_K$, whereas a
nontrivial cyclic cubic conductor is odd and cannot divide sixteen.

Write the minimal polynomial of $H$ as
\[
 z^3-a z^2+(r-1)h z-rh.
\]
The reciprocal identity gives this form with positive integers $a,h$.
The d-number criterion gives $rh\mid a^3$ and
$(rh)^2\mid((r-1)h)^3$. Thus $r^2\mid h$ and $r\mid a$.
The order discriminant $r\operatorname{Norm}(H)=r^2h$ is a square,
so $h=r^2s^2$ and $a=rA$ for positive integers $A,s$. This proves
\eqref{cub:eq:cubic-normalized-polynomial}. For the signed algebra the
same argument applies: $2H$ is an actual formal codegree of the
sphericalization, and the d-number property is unchanged by nonzero
rational scaling.

Put $k=r-1$ and $n=s^2$. The three positive reciprocal roots of
\eqref{cub:eq:cubic-normalized-polynomial} have sum $k$, pairwise-product
sum $A/n$, and product $1/n$. Schur's inequality yields
\[
 4kA\le k^3n+9.
\]
If $k\ge3$ and $4A>k^2n$, the positive integer $4A-k^2n$ is at least
four when $k$ or $s$ is even, and at least three otherwise. This
contradicts the displayed inequality except possibly at $k=3$ with
equality. Equality in Schur for positive numbers makes all three roots
equal, contradicting their cubic degree. Hence $A\le k^2n/4$.
For $k=2$, Schur gives $A\le n+1$. At $A=n+1$ the discriminant is
$n(n-4)=s^2(s^2-4)$, which cannot be a nonzero square. Thus $A\le n$
in this case too. Finally the inverse-square sum is
\[
 \frac{1+k^2-2A/n}{r^2}
 \ge\frac1{r^2}+\frac{k^2}{2r^2},
\]
as claimed.
\end{proof}

\begin{theorem}[Direct cubic reduction]\label{cub:thm:cubic-reduction}
The only possibilities are
\[
 (F,f,g)=(7,3,3),\quad(7,4,4),\quad(9,3,3),
 \qquad \operatorname{Norm}(E)=\operatorname{Norm}(D).
\]
In the conductor-nine case $u=v=1/3$, the ordinary and signed cubic
codegrees agree as multisets, and the category is spherical.
\end{theorem}
\begin{proof}
For a constituent $A$ from the ordinary cubic orbit,
$a_A=1/(2D_i)$ gives $b_A=1/(2f)$ in this measure.
Since $b_A=\rho(F(A))/(2g)$ and $\rho$ is integral-valued,
$g/f\in\ZZ_{>0}$.
Put $a=\operatorname{Tr}(D)$, $b=\operatorname{Tr}(E)$,
$c=\operatorname{Norm}(D)$ and
$k=\operatorname{Norm}(E)/c\in\ZZ_{>0}$, using central
integrality. Casimir trace domination gives $b+g\le a+f$.
For
\[
 H(r)=\frac12-\frac1r+\frac{3}{2r^2},\qquad
 J(r)=1-\frac2r+\frac2{r^2},
\]
Lemma~\ref{cub:lem:cubic-square-norm-schur} and the reciprocal
identities give
\[
 u\ge H(f),\quad u=J(f)-2a/c,\quad v=J(g)-2b/(kc),
 \quad v-u\ge J(g)-J(f).
\]
If $g\ge2f$, substitution in \eqref{cub:eq:cubic-periodic} yields
\[
 \mathcal L\ge\frac32+\frac1f-\frac1{2f^2}-\frac2g+\frac1{g^2}
 \ge\frac32-\frac1{4f^2}\ge\frac{53}{36}>\frac97,
\]
a contradiction. Therefore $g=f$. Now
\[
 \mathcal L\ge\frac32-\frac1f+\frac1{2f^2}.
\]
This increases for $f>1$ and equals $33/25>9/7$ at $f=5$.
Consequently $f=g\in\{3,4\}$.

Lemma~\ref{cub:lem:cubic-square-norm-schur} makes $k$ an integer
square. If $k>1$, then $k\ge4$, and $b\le a$ implies
$v\ge\tfrac34J(f)+\tfrac14u$. Thus
\[
 \mathcal L\ge\frac{33}{16}-\frac{17}{8f}+\frac{17}{16f^2}
 \ge\frac{53}{36}>\frac97,
\]
again impossible. Hence $k=1$.
At $f=3$ the lower bound for $\mathcal L$ is $11/9$, forcing
$F\le9$; at $f=4$ it is $41/32$, forcing $F\le64/9$.
The only cyclic cubic conductors in these ranges are seven and nine,
as follows from primitive Dirichlet characters of order three.
This gives the stated list.

At $(F,f)=(9,3)$ equality holds throughout, so $u=v=1/3$.
Equality of norms and inverse-square sums gives
$\operatorname{Tr}(E)=\operatorname{Tr}(D)$, and the reciprocal
identity identifies the remaining coefficient of the two cubic
polynomials. Their roots therefore agree. Choose a signed character
with codegree $D$ and apply spectral saturation in
Lemma~\ref{core:lem:family-induction} to obtain a spherical section.
\end{proof}

\subsection{An inert-prime obstruction for a rational codegree}

\begin{lemma}\label{cub:lem:rational-prime-inert}
Let $R$ be a commutative self-dual based ring with
$R_{\Q}\cong\Q\times K$, where $K$ is a number field. Suppose the
character given by projection to $\Q$ has formal codegree a prime $p$.
If $L=K\otimes\Q_p$ is a field, the projection of $R\otimes\ZZ_p$
to $L$ is $\mathcal O_L$, and $L/\Q_p$ is totally ramified.
\end{lemma}
\begin{proof}
Let $A=R\otimes\ZZ_p$ and project $A$ to an order $B$ in
$\mathcal O_L$, where $L=K\otimes\Q_p$. If $\rho$ is the rational
character, its idempotent $(1,0)=p^{-1}\sum_i\rho(b_i)b_i$ shows that
$A\cap(\Q_p\times0)=p\ZZ_p\times0$ by looking at the unit coefficient.
Thus the first coordinate modulo $p$ gives a surjection $B\to\mathbb F_p$
and $A=\ZZ_p\times_{\mathbb F_p}B$. The order $B$ is local by
lying-over from its normalization $\mathcal O_L$, so its maximal ideal
$\mathfrak m$ is $B\cap\mathfrak m_L$ and its residue field is
$\mathbb F_p$. If $B=\mathcal O_L$, this gives the claim.

Suppose that $B\ne\mathcal O_L$. The finite $B$-module
$\mathcal O_L/B$ is nonzero and has nonzero socle.
Choose $y\in\mathcal O_L\setminus B$ with $\mathfrak m y\subset B$.
In fact $\mathfrak m y\subset\mathfrak m$ since $y$ is integral.
In $A/pA$ the classes of $(p,0)$ and $(0,py)$ are then both killed by
the maximal ideal, whose inverse image is $p\ZZ_p\times\mathfrak m$.
They are linearly independent over $\mathbb F_p$: a relation would give
$(c,dy)\in A$ after division by $p$, forcing first $p\mid d$ since
$y\notin B$, then $p\mid c$ by the fiber condition (here $py\in\mathfrak m$).
But the coefficient-of-unit pairing on $A/pA$ is associative and perfect.
It identifies the annihilator of the maximal ideal with the dual of the
residue field, of dimension one. This contradiction proves the claim.
The general Gorenstein fiber-product criterion of Nasseh, Sather-Wagstaff,
Takahashi and VandeBogert \cite[Fact~2.6 and Corollary~2.7]{NTV} also
implies the obstruction.
\end{proof}

\begin{corollary}\label{cub:cor:cubic-seven-three-impossible}
There is no commutative self-dual based ring whose rational algebra
is $\Q\times\Q(\zeta_7+\zeta_7^{-1})$ and whose
rational formal codegree is three.
\end{corollary}
\begin{proof}
The polynomial $x^3+x^2-2x-1$ of
$\zeta_7+\zeta_7^{-1}$ has discriminant $49$ and has no root
modulo three: its reduction is $x^3+x^2+x-1$, taking nonzero
values at all three elements of $\mathbb F_3$. Thus three is
unramified and inert in this cubic field. Apply
Lemma~\ref{cub:lem:rational-prime-inert} with $p=3$.
\end{proof}

\subsection{Sphericality at rational codegree four}

\begin{lemma}\label{cub:lem:cubic-four-trace-gap}
In the remaining conductor-seven case, put
$a=\Tr(D)$, $b=\Tr(E)$, and $c=\operatorname{Norm}(D)
=\operatorname{Norm}(E)$. Then
\[
 a-b\in64\ZZ_{\ge0}.
\]
\end{lemma}
\begin{proof}
Let $B$ be the constituent of $I(\mathbf1)$ corresponding to the
rational signed character. Let $C_R$ and $C_T$ be the ordinary and
signed Casimir elements. We show that the ordinary folded and signed
forgetful coefficient vectors of $B$ are respectively
\[
 w=C_R/4,\qquad r=C_T/4.
\]
The signed cubic component of $F(B)$ is its dimension $E/4$.
Its unit coefficient is one, and its sign-object coefficient is zero,
because $B$ belongs to $I(\mathbf1)$ and the supports of $I(\mathbf1)$
and $I(\varepsilon)$ are disjoint. The Frobenius trace formula then
makes its signed rational component one, so its signed class is $C_T/4$.
For an ordinary cubic induction constituent $A_i$, the rational signed
row, already computed in Theorem~\ref{cub:thm:cubic-reduction},
has $S_{B,A_i}=1/8$. Symmetry and the modular character formula give
\[
 \chi_i(F(B))=\frac{S_{A_i,B}}{S_{A_i,0}}
             =\frac{1/8}{1/(2D_i)}=D_i/4.
\]
Its ordinary folded unit coefficient is likewise one, so the ordinary
rational component is one and its folded class is $C_R/4$.

These are two coefficient vectors of the same actual object. Thus
$w_i\ge|r_i|$ and $w_i\equiv r_i\pmod2$ for every $i$. In the
orthonormal based trace, the squared coefficient norm of a Casimir is
the sum of its formal codegrees. Consequently
\[
 a-b=16\bigl(\|w\|^2-\|r\|^2\bigr)\in64\ZZ_{\ge0}.
\]
\end{proof}

\begin{theorem}\label{cub:thm:cubic-four-spherical}
Every category in the conductor-seven, rational-codegree-four case
admits a spherical structure. Its signed and ordinary cubic codegree
polynomials coincide.
\end{theorem}
\begin{proof}
Retain $a,b,c$ from Lemma~\ref{cub:lem:cubic-four-trace-gap}. Here
\[
 u=\frac58-\frac{2a}{c},\qquad
 v=\frac58-\frac{2b}{c},\qquad
 u,v\ge\frac{11}{32}.
\]
The periodic inequality is $u+v\le29/42$. It follows that
\begin{equation}\label{cub:eq:cubic-four-trace-interval}
 \frac{187c}{1344}\le b\le a\le\frac{9c}{64},
 \qquad a-b\le\frac{c}{672}.
\end{equation}
If $a>b$, the trace-gap lemma gives $c\ge43008$.

For either triple of cubic codegrees let $H$ be its largest element
and put $x=1/H$. Since $H$ is at least a third of the trace,
\eqref{cub:eq:cubic-four-trace-interval} gives $0<x<1/1900$.
Let $q_1,q_2$ be the other two reciprocal codegrees. Their sum is
$3/4-x$. Also $u,v\le29/42-11/32=233/672$, so
\[
 q_1^2+q_2^2+x^2\le\frac{191}{672},\qquad
 (q_1-q_2)^2\le\frac1{168}+\frac32x-3x^2.
\]
The last upper bound is less than $(1/12-x)^2$, since the difference is
\[
 \frac1{1008}-\frac53x+4x^2
 >\frac1{1008}-\frac1{1140}>0.
\]
Therefore $1/3<q_1,q_2<5/12$. The two smaller codegrees of each
triple lie in $(12/5,3)$.

Align the largest signed codegree $E_{\max}$ with the ordinary
Frobenius--Perron codegree $D_{\rm FP}$. Central induction makes
$z=E_{\max}/D_{\rm FP}$ an algebraic integer: it gives integrality
against every conjugate of $D$, and one may then apply a field
embedding sending $E$ to $E_{\max}$. Equality of norms makes $z$
a totally positive unit of norm one. Spectral domination gives
$z\le1$ at the chosen embedding. Its other two conjugates are ratios
of numbers in $(12/5,3)$ and are less than $5/4$. Hence
\[
 3\le\Tr(z)<1+2\cdot\frac54=\frac72.
\]
The trace is an integer, so it is three; equality in the
arithmetic--geometric mean inequality forces every conjugate to be
one. Thus $z=1$, contradicting $a>b$.
We conclude $a=b$. Equality of norms and the reciprocal identity
then identifies the two cubic polynomials. Spectral saturation in
Lemma~\ref{core:lem:family-induction} supplies the spherical structure.
\end{proof}

\begin{theorem}\label{cub:thm:cubic-classification}
Every rank-four complex fusion category whose Frobenius--Perron
dimension field has degree three admits a spherical structure.
\end{theorem}
\begin{proof}
Theorem~\ref{cub:thm:cubic-reduction} gives sphericality in conductor
nine. In conductor seven, Corollary~\ref{cub:cor:cubic-seven-three-impossible}
excludes rational codegree three, and Theorem~\ref{cub:thm:cubic-four-spherical}
gives sphericality at codegree four.
\end{proof}

\section{The quartic dimension field}\label{sec:quartic-bound}

\begin{theorem}\label{thm:quartic-uniform-bound}
Let $\C$ have rank four and self-dual simples. If
$D=\FPdim(\C)$ has degree four and $D\ge950$, then $\C$ is spherical.
\end{theorem}

Suppose throughout that $D\ge950$ and that $\C$ has no spherical
structure. By Corollary~\ref{col:cor:rank-four-dimension-field-descent}
the ordinary and signed dimension characters identify their rational
algebras with the same totally real abelian quartic field $K=\Q(D)$.
Index its embeddings by $i=0,1,2,3$, with $0$ the identity.
Choose the signed dimension character so that its codegree $E$ is the
largest signed codegree. Write $D_i,E_i$ for the four conjugates,
including repetitions for $E$. Set
\[
 p_i=D_i^{-1},\qquad q_i=E_i^{-1},\qquad
 d=p_0\le d_*:=1/950,\qquad r=q_0,\qquad
 \alpha=E/D,\quad \alpha_i=p_i/q_i .
\]
Both vectors are positive and sum to one by character orthogonality.
We have $d\le r=\min_i q_i$. In addition, $\alpha_0$ is the least
conjugate of $\alpha$. Indeed the signed and folded dimensions of the
same central column are $E/E_i$ and $D/D_i$
(Proposition~\ref{col:prop:positive-Galois-columns}). The absolute
values of the signed column coefficients are bounded by the ordinary
coefficients, and the absolute signed simple dimensions by the corresponding
FP dimensions. Hence $E/E_i\le D/D_i$. By central integrality and
\eqref{core:discriminants}, $\alpha$ is a totally positive integral
$d$-number of positive square norm $h$. Here the $d$-number property
follows from that of the two codegrees: all conjugates of their ratio
generate the same fractional ideal. In particular
\begin{equation}\label{eq:quartic-products}
 \prod_i q_i=\frac1h\prod_i p_i
 \le \frac{d(1-d)^3}{27}.
\end{equation}
We will rule out $\alpha\ne1$. Equality $E=D$ gives the required
spherical section by Lemma~\ref{core:lem:family-induction}.

\subsection{Moment bounds and the size of the trace}

Write $u=\sum p_i^2$, $v=\sum q_i^2$, and
$K_*=\Q(\sqrt2,\sqrt5)$. The two moment estimates used below are
\begin{equation}\label{eq:quartic-budgets}
 \begin{array}{c|c}
 K\ne K_*&u+v\le4/5\\
 K=K_*&5u+4v\le15/4 .
 \end{array}
\end{equation}
For the first row, Corollary~\ref{core:quartic-short-moments} gives
the assertion if $\sqrt2\notin K$. Otherwise $\sqrt5\notin K$ and
$2u+v\le9/8$, while $u\ge d^2+(1-d)^2/3>13/40$.
The second row is Lemma~\ref{core:lem:sqrt5-periodic-moments}.

\begin{lemma}\label{quart:outside-trace}
If $K\ne K_*$, then $h\in\{1,4\}$, $r<1/150$, and
$\Tr_{K/\Q}\alpha\le40$.
\end{lemma}
\begin{proof}
Trace domination in Lemma~\ref{core:lem:family-trace-norm} gives
$e_3(q)\le e_3(p)/h$, where $e_j$ denotes the elementary symmetric
function. The three coordinates other than $d$ have sum $1-d$, so
$e_3(p)\le (1-d)^2(1+8d)/27$. Four-variable Schur gives
$2v+9e_3(q)\ge1$. If $h\ge9$, then
\[
 u+v\ge\tfrac12+(1-d)^2(17-8d)/54>4/5.
\]
Thus the positive square $h$ equals one or four.

For the remaining estimates put
$c=4/5-d_*^2-(1-d_*)^2/3$. Then $v\le c$.
If every $q_i\ge b:=1/150$, the numbers $x_i=q_i-b$ have sum
$S=1-4b$ and square sum at most $W=c-2b+4b^2$. Homogeneous Schur
gives $9e_3(x)\ge S(S^2-2W)$. For completeness, this four-variable
inequality follows from three-variable Schur by merging the two
smallest variables $y,z$. The difference of the two Schur expressions
is $yz[5(S-y-z)-4(y+z)]\ge0$. Expanding the shifted product gives
\[
 \prod_i q_i\ \ge \frac{bS(S^2-2W)}9+\frac{b^2(S^2-W)}2+b^3S+b^4
 >\frac{d_*(1-d_*)^3}{27}.
\]
The last inequality is direct rational substitution, contradicting
\eqref{eq:quartic-products}. Hence $r<b$. For each $i\ne0$, the other
two coordinates have sum at most $\sqrt{2v}$. Consequently
$q_i\ge1-r-\sqrt{2c}>1/40$, as
$(1-b-1/40)^2>2c$. Finally
$\Tr\alpha= d/r+\sum_{i\ne0}p_i/q_i<1+40$. The trace is integral.
\end{proof}

\begin{lemma}\label{quart:fixed-trace}
If $K=K_*$, then $r<1/80$, $q_i>1/16$ for $i\ne0$, and
$\Tr_{K/\Q}\alpha\le16$.
\end{lemma}
\begin{proof}
For any conjugation $\sigma$ of $K$, the element
$\sigma(\alpha)=(E/\sigma D)/(E/\sigma E)$ belongs to the signed
order $T$. Both displayed quotients belong to $T$ by central induction,
and $E/\sigma E$ has norm one and hence its inverse belongs to $T$.
For a nonidentity involution interchanging $0,i$ and $j,k$, the
unit coefficient of $\alpha\sigma(\alpha)$ is a positive integer
by~\eqref{core:frobenius-trace}. With $U=dp_i$, $V=p_jp_k$, this says
\[
 I=U(E_0+E_i)+V(E_j+E_k)\ge1.
\]
Replacing $E$ by $D$ on the right gives one, since $\sum p_i=1$.
Subtracting the two expressions gives
\[
 I-1=U(\Tr E-\Tr D)+(V-U)(E_j+E_k-D_j-D_k).
\]
In fact $V>U$. Otherwise $\prod p_i=UV\le U^2<d^2$, while
$r^3(1-3r)\le\prod q_i$. This last bound follows by concentrating
all mass above the common lower bound $r$ in one entry. As
$x^3(1-3x)$ increases up to $1/4$, we obtain $r<1/90$.
Now $v\ge r^2+(1-r)^2/3$ and \eqref{eq:quartic-budgets} give
$u<49/100$. Since
$u+2p_jp_k=d^2+p_i^2+(p_j+p_k)^2\ge d^2+(1-d)^2/2$,
we get $V>1/200-d/2>d>U$, a contradiction. Trace domination
therefore gives $E_j+E_k\ge D_j+D_k$ for each pair $j,k\ne0$.

Choose $q_1$ smallest among the three coordinates other than $r$,
put $S=r+q_1$, and write $x=p_1$. The pair bound for $2,3$ gives
$q_2q_3\le(p_2+p_3)(1-S)/4$, and hence
\[
 u\ge d^2+x^2+(1-d-x)^2/2,\qquad
 v\ge r^2+q_1^2+(1-S)^2-(1-d-x)(1-S)/2.
\]
Minimize the resulting quadratic lower bound for $5u+4v$ over real
$x$, and use $rq_1\le S^2/4$. This gives
\[
 5u+4v\ge\frac{21}5-\frac{32S}5+\frac{88S^2}{15}
             -\frac{2d}{3}(2S+3).
\]
If $S\le3/40$, the right side is at least its value at
$(S,d)=(3/40,1/950)$, namely $15/4+3/3800$. Thus $S>3/40$.
By~\eqref{eq:quartic-products}, $r^3(1-3r)<1/(27\cdot950)$, so
$r<1/27$. Suppose $r\ge1/80$. Each remaining coordinate exceeds
$3/40-r$, whence concentrating the excess in one entry gives
\[
 \prod q_i>\psi(r),\qquad
 \psi(r)=r(3/40-r)^2(17/20+r).
\]
On $[1/80,1/27]$ the polynomial is concave
($\psi''=12r^2+21r/5-39/160<0$), and its smaller endpoint value is
$\psi(1/80)=69/1638400>1/(27\cdot950)$. This contradicts
\eqref{eq:quartic-products}, proving $r<1/80$. Now
$q_1>3/40-1/80=1/16$, and $\Tr\alpha<1+16$.
\end{proof}

\subsection{One square inequality for both fields}

For a row $(A,B)$ of \eqref{eq:quartic-budgets} and positive
$\lambda,\mu$, complete the square at each $i\ne0$, using
$p_i=\alpha_i q_i$. Summing gives
\begin{equation}\label{quart:weighted-square}
 Au+Bv\ge 2\lambda(1-d)+2\mu(1-r)
       -\sum_{i=1}^3\frac{(\lambda\alpha_i+\mu)^2}{A\alpha_i^2+B}.
\end{equation}
Only the three conjugates remaining after deleting one least conjugate
enter the correction. We give the full arithmetic test for this
correction so that the quartic argument can be checked separately from
the fusion-ring census.

\begin{lemma}[Finite conjugate estimate]\label{quart:arithmetic}
Let $\alpha\ne1$ be a totally positive integral $d$-number in the
totally real abelian quartic field $K$, of square norm and least conjugate
at most one. Count conjugates with their multiplicities as embeddings
of $K$. For each row of the table, the sum in
\eqref{quart:weighted-square} is less than $M$.
\[
\begin{array}{c|cc|cccc|c}
 &\Tr\alpha\le& h&A&B&\lambda&\mu&M\\\hline
 K\ne K_* &40&1,4&1&1&1/3&1/2&17/20\\
 K=K_* &16&\text{square}&5&4&1&3&41/10
\end{array}
\]
\end{lemma}
\begin{proof}
A rational $\alpha$ would be one. In the first row a quadratic minimal
polynomial is $X^2-tX+n$ with $3\le t\le20$, $n=1,2$, and
$t^2-4n$ positive and nonsquare. If $n=2$, the $d$-number condition
requires $2\mid t$. Each root occurs twice.

For a quartic minimal polynomial $F=X^4-tX^3+sX^2-kX+h$, the
abelian group fixes a partition of the four roots into two pairs.
Let $a,a'$ be their sums, $b,b'$ their products, and put $z=b+b'$.
This is a rational algebraic integer. From the elementary symmetric
functions one obtains
\[
 A_0=(a-a')^2=t^2-4s+4z,\quad B_0=(b-b')^2=z^2-4h,\quad
 (tz-2k)^2=A_0B_0.
\]
Thus the following integer loops cover the first row. Use $4\le t\le40$,
$2\sqrt h\le z\le t^2/4-2\sqrt h$,
$z<s\le3t^2/8$. These bounds follow from AM--GM on the pairs and on
the elementary symmetric sums. Test $A_0,B_0\ge0$, $A_0B_0$ square
and $k=(tz\pm\sqrt{A_0B_0})/2$ integral in $[1,t^3/16]$.
When $h=4$ retain only $2\mid t,s$ and $4\mid k$.
In general, if $c_j$ is the $j$th signed elementary symmetric
coefficient, the $d$-number property gives $h^j\mid c_j^4$:
the four conjugates generate the same ideal.

The discriminants of the two quadratic factors over their pair field
have sum $t^2-2s-2z$ and product
$J=s^2+2sz-3z^2-4tk+16h$. Retain positive sum and product.
Remove reducible polynomials by testing monic linear and quadratic
integer factors with positive roots and constant term dividing $h$.
The pair field is $\Q(\sqrt\delta)$, with $\delta=A_0$ if $A_0>0$ and
$\delta=B_0$ otherwise. Irreducibility requires $\delta$ nonsquare.
Since $K$ splits both quadratic factors and is quadratic over the pair
field, their discriminants represent the same square class there.
Hence $J$ or $J\delta$ must be an integer square. Retain either
possibility and count equal polynomials only once.

For the second row put $\phi=(1+\sqrt5)/2$.
The integral basis $1,\sqrt2,\phi,\sqrt2\phi$ follows from the
coprime quadratic discriminants $8,5$. Write
\[
 4\alpha=t+b\sqrt2+c\sqrt5+e\sqrt{10}.
\]
Then $2\le t\le16$ is even, $c,e$ are even, and
$t\equiv c$, $b\equiv e\pmod4$. Total positivity gives
$2b^2,5c^2,10e^2<t^2$ by adding and subtracting conjugates.
Enumerate these integers, and form the polynomial of the four conjugates
by putting $y=4X-t$ in
\[
 4^4F(X)=(y^2+5c^2-2b^2-10e^2)^2-5(2cy+4be)^2.
\]
Keep alternating nonzero coefficient signs (all conjugates are real),
square norm $h$, $h^j\mid c_j^4$, and least root at most one.
Remove rational constants. There are two distinct roots precisely
when just one of $b,c,e$ is nonzero.

Here are the bounds obtained from these loops. For each distinct root
use its adjacent grid points of mesh $2^{-16}$, duplicating the
intervals when the degree is two. For ordered intervals $[a_i,b_i]$,
a rational upper bound on the desired sum is
$\sum_{i=1}^3(\lambda b_i+\mu)^2/(Aa_i^2+B)$. A grid interval for the
$j$th root is found by bisection on $[0,t]$. The number of roots
larger than a grid point $x$ equals the number of sign changes in the
nonzero coefficients of $F(X+x)$ by Descartes' rule, which is exact
when all roots are real. This specifies the entire finite check.
\[
\begin{array}{c|rr|c|l}
 &\deg2&\deg4&\text{upper bound}&
 \text{maximizing polynomial}\\\hline
 K\ne K_*&27&291&841/1000&X^4-16X^3+20X^2-8X+1\\
 K=K_*&5&3&4068/1000&X^2-3X+1
\end{array}
\]
The counts refer to the arithmetic candidates tested, before any
fusion-ring conditions. The last column gives a minimal polynomial attaining the largest rational interval bound. Both upper bounds are below the corresponding
$M$. The accompanying script \texttt{quartic\_bounds.py} implements
these loops with integers and rational endpoints.
\end{proof}

\begin{proof}[Proof of Theorem~\ref{thm:quartic-uniform-bound}]
If $\alpha\ne1$, Lemmas~\ref{quart:outside-trace}
and~\ref{quart:fixed-trace} allow us to apply
Lemma~\ref{quart:arithmetic} in~\eqref{quart:weighted-square}.
In the two cases, respectively, this yields
\[
 u+v>\frac53-\frac2{3\cdot950}-\frac1{150}-\frac{17}{20}>\frac45,
 \qquad
 5u+4v>8-\frac2{950}-\frac6{80}-\frac{41}{10}>\frac{15}4.
\]
Each inequality contradicts its moment bound. Thus $E=D$, and
spectral saturation gives a spherical section.
\end{proof}

\section{The dimension bound in the nonspherical case}\label{sec:global-bound-proof}
\begin{proof}[Proof of Theorem~\ref{thm:uniform-global}]
Suppose $\C$ has no spherical structure. By
Theorem~\ref{ns:thm:ns-spherical} every simple is self-dual.
If $D$ were rational, it would be integral, and weak integrality gives
a spherical structure by \cite[Proposition~8.24]{ENO}.
The Frobenius--Perron character field is a factor of the rank-four
rational algebra and contains $\Q(D)$. If $D$ has degree three,
so does the character field, and
Theorem~\ref{cub:thm:cubic-classification} gives sphericality.

Only degrees two and four remain. In either case
Corollary~\ref{col:cor:rank-four-dimension-field-descent} identifies
the signed dimension field with $\Q(D)$. In degree two,
Theorem~\ref{qs:short-orbits} gives a product of two copies of this
field for the signed algebra, with complementary codegree product at
least $100$. Theorem~\ref{quad:opt:unconditional} then gives $D<1000$.
In degree four Theorem~\ref{thm:quartic-uniform-bound} gives $D<950$.
Finally Lemma~\ref{core:lem:coefficient-bound} gives
$N_{XY}^{Z}\le d_X<\sqrt D<32$. As these coefficients are integers,
all are at most $31$.
\end{proof}
\section{Completeness of the bounded census}\label{app:ordinary-census}
The ordinary coefficients lie in $[0,M]\cap\ZZ$, with $M=31$.
We keep rings of global Frobenius--Perron dimension strictly below
$\Lambda=1000$. Dong, Palcoux and Plessis used an ordinary census in their pivotal
classification \cite{DPP}. Here a common linear search gives the ordinary
census and the local signed lifts needed for descent.

\subsection{Presentations and elimination}\label{finite:sec:census}
Use the row convention $(N_u)_{vw}=N_{uv}^w$ and put
$B=\sum_u N_uN_u^{\mathsf T}$, the ordinary Casimir.
Every rank-four fusion ring is commutative by
Lemma~\ref{core:lem:rank-four-algebra}. For a self-dual basis $1,x,y,z$, reciprocity gives
\[
N_x=\begin{pmatrix}0&1&0&0\\1&a&b&c\\0&b&d&e\\0&c&e&f\end{pmatrix},\quad
N_y=\begin{pmatrix}0&0&1&0\\0&b&d&e\\1&d&g&h\\0&e&h&i\end{pmatrix},\quad
N_z=\begin{pmatrix}0&0&0&1\\0&c&e&f\\0&e&h&i\\1&f&i&j\end{pmatrix}.
\]
These formulas also define signed tables, when arbitrary integer coefficients are allowed.
Associativity is equivalent to the following six equations:
\begin{align}
 ad-b^2+bg+ch-d^2-e^2+1&=0,\label{finite:eq:census1}\\
 ae-bc+bh+ci-de-ef&=0,\label{finite:eq:census2}\\
 be-cd+dh-eg+ei-fh&=0,\label{finite:eq:census3}\\
 af+bi-c^2+cj-e^2-f^2+1&=0,\label{finite:eq:census4}\\
 bf-ce+di-eh+ej-fi&=0,\label{finite:eq:census5}\\
 df-e^2+gi-h^2+hj-i^2+1&=0.\label{finite:eq:census6}
\end{align}
The six equations are obtained by comparing the three pairwise matrix commutators.

The remaining duality has a single pair. For duality $(0)(1)(23)$ the parameters are $(a,b,c,d,e,f)$,
and
\[
 N_1=\begin{pmatrix}0&1&0&0\\1&a&b&b\\0&b&c&d\\0&b&d&c\end{pmatrix},\quad
 N_2=\begin{pmatrix}0&0&1&0\\0&b&c&d\\0&d&e&f\\1&c&e&e\end{pmatrix},\quad
 N_3=N_2^{\mathsf T}.
\]
The displayed matrices commute exactly when
\begin{align}
 ac-b^2+2be-c^2-d^2+1&=0,\notag\\
 ad-b^2+be+bf-2cd&=0,\label{census:eq:ns-assoc}\\
 b(c-d)+d(e-f)&=0,\notag\\
 -c^2+d^2-e^2+f^2-1&=0.\notag
\end{align}
In both cases commutativity of the displayed matrices suffices for
associativity of the proposed multiplication. Indeed, set
$E=N_uN_v-\sum_w N_{uv}^{w}N_w$. Its zeroth row vanishes, and it
commutes with every $N_t$. Its row $t$ is the zeroth row of
$N_tE=EN_t$, hence also vanishes. Thus $E=0$, proving the full
regular multiplication identities. The unit, duality, integrality,
nonnegativity, and reciprocity axioms are already built into the
presentations.

\subsection{A single bounded linear search}\label{sec:linear-search}
For any finite integer domain $V$, write
$\mathcal L(k,r;V)=\{u\in V:ku=r\}$. For $k\ne0$ this is the singleton
$\{r/k\}$ if $r/k\in V$ and is empty otherwise. For $k=0$ it is
$V$ if $r=0$ and empty otherwise. This convention includes every free
variable on a degenerate branch.

For a self-dual table use $V=\{0,\ldots,M\}$. Permutations of the
three nonunit basis elements are transitive on the six coordinates
$b,c,d,f,h,i$, so we may put a minimal coordinate at $b$.
Enumerate $a,b,c,d,e,f,h\in V$ with $b\le c,d,f,h$. The first,
second and fourth equations give, in order,
\[
\begin{aligned}
g&\in\mathcal L(b,b^2+d^2+e^2-1-ad-ch;V),\\
i&\in\mathcal L(c,bc+e(d+f-a)-bh;V), \qquad i\ge b,\\
j&\in\mathcal L(c,c^2+e^2+f^2-1-af-bi;V).
\end{aligned}
\]
Test the other three equations. Keep the lexicographically least tuple
among the permutations for which the minimum is still at $b$.
The only early dimension cutoff uses
\[
B_{11}=2+a^2+2b^2+2c^2+d^2+2e^2+f^2\le D .
\]
Thus any choice with $B_{11}\ge\Lambda$ can be skipped before enumerating
$h$. The same three linear equations apply to signed lifts when $V$ is
replaced in each coordinate by its permitted parity interval.

For a dual pair enumerate $c,d,e,f\in V$ and test the last equation of
\eqref{census:eq:ns-assoc}. For each $b\in V$ satisfying the third,
take $a\in\mathcal L(c,b^2-2be+c^2+d^2-1;V)$ and test the second.
Exchanging the dual pair leaves all six parameters unchanged. Distinct
tuples therefore give distinct based rings.

\subsection{The dimension test and the coverage conclusion}
The positive dimension vector is an eigenvector of the symmetric
nonnegative Casimir matrix $B$, with eigenvalue $D=\FPdim(R)$.
It follows that $D$ is its largest eigenvalue. Therefore
\begin{equation}\label{census:eq:pd}
 D<\Lambda\quad\Longleftrightarrow\quad
 \Lambda I-B\text{ is positive definite}
 \quad\Longleftrightarrow\quad
 \det(\Lambda I-B)_{[1,k]}>0\ (1\leq k\leq4).
\end{equation}
Here $[1,k]$ denotes the leading principal $k$-by-$k$ submatrix.
The final equivalence is Sylvester's criterion. All determinants are
integers; a zero determinant is rejected, so equality at the boundary
is never admitted. The C++ enumeration uses fraction-free elimination,
and the Python tensor checker recomputes the same quantities with
arbitrary-precision integers.

\begin{proposition}\label{census:prop:coverage}
The preceding procedure with $M=31$ and $\Lambda=1000$ emits exactly
one representative of every rank-four fusion ring with $D<1000$.
\end{proposition}
\begin{proof}
Every rank-four fusion ring has one of the two complete presentations.
All solutions of their associativity equations in the parameter box
occur in the enumerated branches, and the isomorphism reduction retains
exactly one representative. The acceptance condition
\eqref{census:eq:pd} is equivalent to the desired strict dimension
bound. Conversely, every ring below that bound has $\mu<\sqrt D$
by Lemma~\ref{core:lem:coefficient-bound}, hence $\mu\leq31$, and therefore
lies in the searched box. This proves both directions.
\end{proof}

The executable output has $1642$ self-dual classes and
$50$ classes with a dual pair, for a total of $1692$.
As an additional implementation check, removing both dimension-pruning
steps gives $16065$ classes at multiplicity at most $31$; an independently
implemented strict dimension test retains precisely the same $1692$
parameter tuples. A full small-box test examines all
$3^{10}+3^6=59778$ assignments at multiplicity at most two and recovers
exactly the same $27$ classes. The unpruned experiment shares the
linear search. The small-box test uses the same
tensor presentations and canonicalization.

\section{Sphericality from the bounded census}\label{sec:finite-descent}

Suppose that a rank-four category has no spherical structure.
Theorem~\ref{thm:uniform-global} puts its ordinary ring $R$ in the census.
Write $P_R$ for the characteristic polynomial of its Casimir matrix and
$P_T$ for the corresponding polynomial of its signed quotient. Their roots,
with multiplicity, are the ordinary and signed formal codegrees.
The finite argument below uses only the coefficients of these polynomials
and the six associativity equations. No sign choice for the spherical
dimension character itself is required.

\begin{lemma}[Two filters]\label{lem:descent-filters}
For such a category, the following conditions hold.
\begin{enumerate}
\item Either $P_R$ is irreducible quartic, or it has only linear and
quadratic factors with exactly one nontrivial quadratic splitting field
among them. Each irreducible factor is the minimal polynomial of a
$d$-number.
\item Let $\mathcal F$ be the set of monic irreducible factors of $P_RP_T$.
For at least one irreducible factor $e$ of $P_T$,
\[
 |f(0)|^{\deg e}\ \bigm|\ |e(0)|^{\deg f}\qquad(f\in\mathcal F).
\]
\end{enumerate}
\end{lemma}
\begin{proof}
The proof of Theorem~\ref{thm:uniform-global} leaves two possibilities for
$\Q(D)$ and identifies it with the signed dimension field $L$. In degree
four the Casimir eigenvalue $D$ has four distinct conjugates and these exhaust
$P_R$. In degree two, all ordinary codegrees lie in $L$ by
Lemma~\ref{core:lem:family-induction}, including the irrational $D$.
The $d$-number assertion is Ostrik's theorem \cite{OstrikCodegrees}.

For the second assertion take $e$ to be the minimal polynomial of the signed
global codegree $E$. Central induction makes $E/h$ an algebraic integer for
every root $h$ represented in $\mathcal F$. Take absolute norms in a number
field $K$ containing $E$ and $h$. At each prime the resulting inequality is
\[
 \frac{v_p(|e(0)|)}{\deg e}\ \geq\
 \frac{v_p(|f(0)|)}{\deg f},
\]
after cancellation of $[K:\Q]$. This gives the divisibility.
\end{proof}

\begin{proposition}[Finite reduction]\label{prop:two-polys}
The associativity, parity, domination and norm conditions below leave
one possible ordinary tuple and two possible signed Casimir polynomials:
\[
\begin{aligned}
 N&=(2,2,2,2,3,4,2,4,4,6),\\
 P_R(t)&=t^4-288t^3+2880t^2-9216t+9216,\\
 P_{T,1}(t)&=t^4-96t^3+1728t^2-9216t+9216,\\
 P_{T,2}(t)&=(t^2-72t+96)(t^2-24t+96).
\end{aligned}
\]
\end{proposition}
\begin{proof}[Exact enumeration]
Start with the self-dual census. A nontrivial invertible gives sphericality
by Theorem~\ref{inv:thm:rankfour-invertible-spherical}. A nonunit
basis object with $X^2=1+nX$ for $n\ge2$ is excluded by
Proposition~\ref{core:prop:rank-two}. For the remaining tensors factor $P_R$
over $\Q$ and apply Lemma~\ref{lem:descent-filters}(1). The $d$-number
test is the coefficient divisibility in Proposition~\ref{finite:prop:standard}.
For quadratic factors compare the squarefree parts of their discriminants.

For each surviving ten-tuple $N=(n_1,\ldots,n_{10})$, the $k$th signed
entry ranges over $-n_k,-n_k+2,\ldots,n_k$. In the order of
\eqref{eq:tuple-convention}, enumerate $a,b,c,d,e,f,h$ in those ranges.
One may take $a,b,c\ge0$ since the entries $T_{111},T_{112},T_{113}$
transform by the independent signs of $X,Y,Z$. Use the three linear
equations in Section~\ref{sec:linear-search} to determine $g,i,j$
in their respective parity intervals, keeping a whole interval when
the corresponding equation is $0=0$. Check the other three equations.
If every absolute entry
equals its ordinary counterpart, Lemma~\ref{trace:lem:absolute-rigidity}
gives a spherical section, so discard it. For each ordinary ring, collect
the distinct remaining $P_T$, computed as
$\det(tI-\sum_i T_iT_i^{\mathsf T})$, and apply
Lemma~\ref{lem:descent-filters}(2). The following counts record the
finite calculation; the last two count pairs $(R,P_T)$.
\begin{center}\small
\begin{tabular}{@{}lr@{}}\toprule
Condition just imposed & Remaining \\ \midrule
Self-dual ordinary rings of $D<1000$ & 1642\\
No nontrivial invertible & 1259\\
No forbidden rank-two subring & 905\\
Ordinary formal codegrees are $d$-numbers & 71\\
Ordinary codegree field test & 41\\ \midrule
Strict signed tuples in the chosen signs, counted over the 41 rings & 1798\\
Distinct pairs $(R,P_T)$ & 623\\
Joint norm divisibility & 2\\ \bottomrule
\end{tabular}
\end{center}
The two pairs are the displayed output. All steps are integer operations or
factorizations of polynomials of degree at most four over $\Q$. The script
\texttt{verify\_descent.py} implements the enumerations just described.
\end{proof}

\begin{theorem}\label{thm:universal-spherical}
Every complex fusion category of rank four admits a spherical structure.
\end{theorem}
\begin{proof}
For the single ordinary candidate, $P_R$ is irreducible, so $[\Q(D):\Q]=4$.
Corollary~\ref{col:cor:rank-four-dimension-field-descent} now identifies the
signed dimension field $L$ with $\Q(D)$. Thus
Proposition~\ref{core:squared-moments} requires $u+v\le7/8$ for either
of the two proposed signed polynomials.

For any quartic $t^4+At^3+Bt^2+Ct+H$ with nonzero roots $f_i$,
Vieta's formulas give $\sum_i f_i^{-2}=(C/H)^2-2B/H$.
The ordinary polynomial gives $u=3/8$. For $P_{T,1}$ and $P_{T,2}$
the respective values of $v$ are $5/8$ and $7/12$. Thus $u+v$ is
$1$ or $23/24$, both strictly greater than $7/8$.
\end{proof}

\section{The fifteen rings and sharp bounds}\label{finite:sec:catalog}
\begin{corollary}[Rank-four classification]\label{finite:cor:classification}
The rings in Table~\ref{finite:tab:models} are exactly the Grothendieck
rings of complex fusion categories of rank four, up to based-ring
isomorphism. Each ring has a unitary realization.
\end{corollary}
\begin{proof}
The preceding descent gives a spherical structure on any rank-four category
under consideration. In particular it is pivotal. The classification of
Dong, Palcoux and Plessis \cite[Theorem~1.1 and Table~1]{DPP} gives exactly fifteen based rings
for complex pivotal categories of rank four. The fifteen rows below are
realized by the indicated categories and are pairwise nonisomorphic, as
one checks from their multiplication tensors by the six permutations of
the nonunit basis elements.
\end{proof}
% Model rows checked by supplementary/models/check_models.py.
\begin{table}[htbp]
\centering\small
\begin{tabular}{@{}clcl@{}}
\toprule
$\mu$ & A unitary realization & $D$ & Reference \\
\midrule
$1$ & $\mathrm{Vec}_{C_2\times C_2}$ & $4$ & \cite[\S3.4]{LPR} \\
$1$ & $\mathrm{Fib}\boxtimes\mathrm{Vec}_{C_2}$ & $5+\sqrt{5}$ & \cite[\S3.4]{LPR} \\
$1$ & $\operatorname{Rep}(D_{10})$ & $10$ & \cite[\S3.4]{LPR} \\
$1$ & $\operatorname{Ad}(\mathrm{SU}(2)_6)$ & $8+4\sqrt{2}$ & \cite[Theorem 3.1]{EIP} \\
$1$ & $\mathrm{Fib}\boxtimes\mathrm{Fib}$ & $(15+5\sqrt{5})/2$ & \cite[\S3.4]{LPR} \\
$1$ & $\operatorname{Ad}(\mathrm{SU}(2)_7)$ & $\delta_7$ & \cite[\S3.4]{LPR} \\
$1$ & $\mathrm{Vec}_{C_4}$ & $4$ & \cite[\S3.4]{LPR} \\
$1$ & Tambara--Yamagami on $C_3$ & $6$ & \cite{TambaraYamagami} \\
$1$ & Principal even part of $S'$ & $8+4\sqrt{2}$ & \cite[Theorem 1]{LMP} \\
$2$ & Principal even part of $2D2$ & $20+8\sqrt{5}$ & \cite[Theorem 3.1]{EIP} \\
$2$ & Haagerup category $\mathcal H_1$ & $(39+9\sqrt{13})/2$ & \cite[Table 2]{GrossmanSnyder} \\
$2$ & $\operatorname{Rep}(A_4)$ & $12$ & \cite[Theorem 1.1]{Larson} \\
$3$ & $\operatorname{Ad}(E_8)$ & $\delta_E$ & \cite{IzumiE8,EvansJones} \\
$3$ & Near-group type $(C_3,3)$ & $(21+3\sqrt{21})/2$ & \cite[Example 9.4]{Izumi} \\
$6$ & Near-group type $(C_3,6)$ & $24+12\sqrt{3}$ & \cite[Theorem 10.19]{Izumi} \\
\bottomrule
\end{tabular}
\caption{One unitary realization of each surviving based ring, with its
multiplicity $\mu$ and global Frobenius--Perron dimension $D$.
Here $D_{10}$ has order ten, $S'$ has index $3+2\sqrt2$, and the
near-group category of type $(C_3,3)$ is the Izumi--Xu category.}
\label{finite:tab:models}
\smallskip
\begin{minipage}{\textwidth}\footnotesize
The numbers $\delta_7$ and $\delta_E$ are the largest real roots of
$x^3-27x^2+162x-243$ and
$x^4-45x^3+375x^2-1125x+1125$, respectively.
The supplementary model check verifies these polynomials, the positive
dimension characters, and all tensor identifications exactly.
\end{minipage}
\end{table}

Table~\ref{finite:tab:models} supplies one unitary realization and a
construction reference for every surviving based ring. Here
$\operatorname{Ad}(\mathrm{SU}(2)_k)$ denotes the adjoint subcategory,
$\mathcal H_1$ is the rank-four Haagerup even part, and a near-group ring
of type $(C_3,m)$ satisfies
$\rho^2=\sum_{g\in C_3}g+m\rho$.
The six non-self-dual models are
$\mathrm{Vec}_{C_4}$, the Tambara--Yamagami category on $C_3$,
the principal even part of $S'$, $\operatorname{Rep}(A_4)$, and the
two near-group categories of parameters three and six.

Asaeda and Haagerup constructed the index-$(5+\sqrt{13})/2$ subfactor
\cite{AsaedaHaagerup}. Grossman and Snyder give the table of its rank-four
even part $\mathcal H_1$ \cite[Table~2]{GrossmanSnyder}. The $2D2$ subfactor is due to
Morrison--Penneys \cite[Theorem~1.1]{MorrisonPenneys}. The $E_8$
construction is due to Izumi \cite{IzumiE8}; its adjoint fusion
rules are displayed in \cite[Section~3.3]{EvansJones}. The
$(C_3,6)$ example was first found by Zhengwei Liu and Noah Snyder,
as credited in \cite[Section~10.6]{Izumi}.

The model check reconstructs each tensor from the group
operations, the product or truncated $\mathrm{SU}(2)$ rules, or the
cited exceptional fusion table. It identifies the corresponding ordinary
census record and checks the positive character, the multiplicity and
the global dimension. The cited constructions supply the unitary categories.
\begin{corollary}[Sharp bounds]\label{finite:cor:sharp-bounds}
Every rank-four complex fusion category satisfies
\[
\operatorname{FPdim}(\C)\leq24+12\sqrt3,
\qquad \max_{i,j,k}N_{ij}^k\leq6.
\]
Both bounds are attained, and equality in either bound holds precisely
for the near-group ring $\operatorname{NG}(C_3,6)$.
\end{corollary}
\begin{proof}
Apply the sharp pivotal bound of Dong, Palcoux and Plessis
\cite[Corollary~5.4]{DPP} after Theorem~\ref{thm:universal-spherical}.
For the equality case the noninvertible simple has dimension
$3+2\sqrt3$, so $D=3+(3+2\sqrt3)^2=24+12\sqrt3$.
\end{proof}

\paragraph{Reproduction.}
From the \texttt{verification} directory run \texttt{python3 verify.py}.
This command regenerates the ordinary census, tests every signed candidate,
checks the quadratic and quartic estimates, and reconstructs the fifteen
model tensors. It requires Python, SymPy and a C++17 compiler.
\begingroup\small
\begin{thebibliography}{99}
\bibitem{AsaedaHaagerup}
M.~Asaeda and U.~Haagerup,
\emph{Exotic subfactors of finite depth with Jones indices
$(5+\sqrt{13})/2$ and $(5+\sqrt{17})/2$},
Comm. Math. Phys. \textbf{202} (1999), 1--63.
\href{https://arxiv.org/abs/math/9803044}{arXiv:math/9803044}.

\bibitem{Bartlett}
B. Bartlett, \emph{Fusion categories via string diagrams},
Commun. Contemp. Math. \textbf{18} (2016), no.~5, 1550080, 39 pp.
\href{https://doi.org/10.1142/S0219199715500807}{doi:10.1142/S0219199715500807}.

\bibitem{Bruillard}
P.~Bruillard, \emph{Rank 4 premodular categories},
with an appendix by P.~Bruillard, C.~Galindo, S.-H.~Ng, J.~Plavnik,
E.~Rowell and Z.~Wang, New York J. Math. \textbf{22} (2016), 775--800.
\href{https://nyjm.albany.edu/j/2016/22-36p.pdf}{Published article}.

\bibitem{DCW}
J.~Dong, G.~Chen and Z.~Wang, \emph{Fusion categories
containing a fusion subcategory with maximal rank}, J. Algebra
\textbf{604} (2022), 107--127.
\href{https://doi.org/10.1016/j.jalgebra.2022.02.020}{doi:10.1016/j.jalgebra.2022.02.020}.

\bibitem{DLN}
C. Dong, X. Lin, and S.-H. Ng, \emph{Congruence property in conformal
field theory}, Algebra Number Theory \textbf{9} (2015), no.~9, 2121--2166.
\href{https://doi.org/10.2140/ant.2015.9.2121}{doi:10.2140/ant.2015.9.2121}.

\bibitem{DongFour}
J.~Dong, \emph{On Grothedieck rings of rank $4$
self-dual fusion categories}, preprint (2025).
\href{https://arxiv.org/abs/2505.19153v2}{arXiv:2505.19153v2}.

\bibitem{DPP}
J. Dong, S. Palcoux and A. Plessis,
\emph{Classifying Grothendieck rings of pivotal fusion categories of rank four},
preprint, 2026.
\href{https://arxiv.org/abs/2609.33272v2}{arXiv:2609.33272v2}.

\bibitem{EGNO}
P. Etingof, S. Gelaki, D. Nikshych and V. Ostrik,
\emph{Tensor Categories}, Mathematical Surveys and Monographs \textbf{205},
American Mathematical Society, 2015.

\bibitem{EIP}
C.~Edie-Michell, M.~Izumi and D.~Penneys,
\emph{Classification of $\ZZ/2\ZZ$-quadratic unitary fusion categories},
with an appendix by R.~Johnson, S.-H.~Ng, D.~Penneys, J.~Roat,
M.~Titsworth and H.~Tucker,
Quantum Topol. \textbf{16} (2025), 493--568.
\href{https://arxiv.org/abs/2108.01564v2}{arXiv:2108.01564v2}.

\bibitem{ENO}
P. Etingof, D. Nikshych, and V. Ostrik, \emph{On fusion categories},
Ann. of Math. (2) \textbf{162} (2005), no.~2, 581--642.
\href{https://doi.org/10.4007/annals.2005.162.581}{doi:10.4007/annals.2005.162.581}.

\bibitem{EvansJones}
D. E. Evans and C. Jones,
\emph{Quantum symmetries of noncommutative tori},
Comm. Math. Phys. \textbf{406} (2025), article 167.
\href{https://doi.org/10.1007/s00220-025-05337-y}{doi:10.1007/s00220-025-05337-y}.

\bibitem{GannonSchopieray}
T.~Gannon and A.~Schopieray,
\emph{Algebraic number fields generated by dimensions in fusion rings},
Commun. Number Theory Phys. \textbf{18} (2024), 705--743.
\href{https://doi.org/10.4310/CNTP.241202233757}{doi:10.4310/CNTP.241202233757}.

\bibitem{GrossmanSnyder}
P.~Grossman and N.~Snyder,
\emph{Quantum subgroups of the Haagerup fusion categories},
Comm. Math. Phys. \textbf{311} (2012), 617--643.
\href{https://doi.org/10.1007/s00220-012-1427-x}{doi:10.1007/s00220-012-1427-x}.

\bibitem{HaggeHong}
T.~J.~Hagge and S.-M.~Hong, \emph{Some non-braided
fusion categories of rank three}, Commun. Contemp. Math. \textbf{11}
(2009), 615--637.
\href{https://doi.org/10.1142/S0219199709003521}{doi:10.1142/S0219199709003521}.

\bibitem{Izumi}
M.~Izumi, \emph{A Cuntz algebra approach to the classification of near-group categories},
in \emph{Proceedings of the 2014 Maui and 2015 Qinhuangdao Conferences
in Honour of Vaughan F. R. Jones' 60th Birthday},
Proc. Centre Math. Appl. Austral. Nat. Univ. \textbf{46} (2017), 222--343.
\href{https://maths.anu.edu.au/files/8\%2520Izumi_2.pdf}{Published chapter}.

\bibitem{IzumiE8}
M. Izumi,
\emph{On flatness of the Coxeter graph $E_8$},
Pacific J. Math. \textbf{166} (1994), 305--327.

\bibitem{Larson}
H. K. Larson, \emph{Pseudo-unitary non-self-dual fusion categories of rank 4},
J. Algebra \textbf{415} (2014), 184--213.
\href{https://doi.org/10.1016/j.jalgebra.2014.05.032}{doi:10.1016/j.jalgebra.2014.05.032}.

\bibitem{LMP}
Z. Liu, S. Morrison and D. Penneys,
\emph{1-supertransitive subfactors with index at most $6\frac15$},
Comm. Math. Phys. \textbf{334} (2015), 889--922.
\href{https://doi.org/10.1007/s00220-014-2160-4}{doi:10.1007/s00220-014-2160-4}.

\bibitem{LPR}
Z.~Liu, S.~Palcoux and Y.~Ren,
\emph{Classification of Grothendieck rings of complex fusion categories
of multiplicity one up to rank six},
Lett. Math. Phys. \textbf{112} (2022), article~54.
\href{https://arxiv.org/abs/2010.10264v5}{arXiv:2010.10264v5}.

\bibitem{MorrisonPenneys}
S.~Morrison and D.~Penneys,
\emph{2-supertransitive subfactors at index $3+\sqrt5$},
J. Funct. Anal. \textbf{269} (2015), no.~9, 2845--2870.
\href{https://doi.org/10.1016/j.jfa.2015.06.023}{doi:10.1016/j.jfa.2015.06.023}.

\bibitem{NS}
S.-H. Ng and P. Schauenburg, \emph{Frobenius--Schur indicators and exponents of spherical categories},
Adv. Math. \textbf{211} (2007), no.~1, 34--71.
\href{https://doi.org/10.1016/j.aim.2006.07.017}{doi:10.1016/j.aim.2006.07.017}.

\bibitem{NSHigher}
S.-H. Ng and P. Schauenburg, \emph{Higher Frobenius--Schur indicators for
pivotal categories}, in \emph{Hopf Algebras and Generalizations},
Contemp. Math. \textbf{441}, Amer. Math. Soc., Providence, RI, 2007, 63--90.
\href{https://doi.org/10.1090/conm/441/08500}{doi:10.1090/conm/441/08500}.

\bibitem{NTV}
S. Nasseh, S. Sather-Wagstaff, R. Takahashi, and K. VandeBogert,
\emph{Applications and homological properties of local rings with
decomposable maximal ideals}, J. Pure Appl. Algebra \textbf{223} (2019),
no.~3, 1272--1287.
\href{https://doi.org/10.1016/j.jpaa.2018.06.006}{doi:10.1016/j.jpaa.2018.06.006}.

\bibitem{Ostrik3}
V. Ostrik, \emph{Pivotal fusion categories of rank 3}, with an appendix
written jointly with D. Nikshych, Mosc. Math. J. \textbf{15} (2015),
no.~2, 373--396.
\href{https://doi.org/10.17323/1609-4514-2015-15-2-373-396}{doi:10.17323/1609-4514-2015-15-2-373-396}.

\bibitem{OstrikCodegrees}
V. Ostrik, \emph{On formal codegrees of fusion categories},
Math. Res. Lett. \textbf{16} (2009), no.~5, 895--901.
\href{https://doi.org/10.4310/MRL.2009.v16.n5.a11}{doi:10.4310/MRL.2009.v16.n5.a11}.

\bibitem{OstrikRankTwo}
V. Ostrik, \emph{Fusion categories of rank 2}, with an appendix by P.~Etingof,
Math. Res. Lett. \textbf{10} (2003), no.~2--3, 177--183.
\href{https://arxiv.org/abs/math/0203255}{arXiv:math/0203255}.\par

\bibitem{RSW}
E.~C. Rowell, R. Stong, and Z. Wang, \emph{On classification of modular
 tensor categories}, Commun. Math. Phys. \textbf{292} (2009), no.~2, 343--389.
\href{https://doi.org/10.1007/s00220-009-0908-z}{doi:10.1007/s00220-009-0908-z}.

\bibitem{Schopieray}
A. Schopieray,
\emph{Categorification of integral group rings extended by one dimension},
J. Lond. Math. Soc. (2) \textbf{108} (2023), 1617--1641;
\href{https://doi.org/10.1112/jlms.12789}{doi:10.1112/jlms.12789}.

\bibitem{SchopierayNorm}
A.~Schopieray, \emph{Norm, trace, and formal
codegrees of fusion categories}, J. Algebra \textbf{568} (2021), 362--385.
\href{https://doi.org/10.1016/j.jalgebra.2020.10.014}{doi:10.1016/j.jalgebra.2020.10.014}.

\bibitem{TambaraYamagami}
D. Tambara and S. Yamagami,
\emph{Tensor categories with fusion rules of self-duality for finite abelian groups},
J. Algebra \textbf{209} (1998), 692--707.
\href{https://doi.org/10.1006/jabr.1998.7558}{doi:10.1006/jabr.1998.7558}.

\bibitem{VS}
G.~Vercleyen and J.~K.~Slingerland,
\emph{On low rank fusion rings}, J. Math. Phys. \textbf{64} (2023), 091703.
\href{https://doi.org/10.1063/5.0148848}{doi:10.1063/5.0148848}.

\bibitem{VSData}
G.~Vercleyen and J.~Slingerland, \emph{On low-rank
multiplicity-free complex fusion categories}, Comm. Math. Phys. \textbf{407}
(2026), article 231.
\href{https://doi.org/10.1007/s00220-026-05705-2}{doi:10.1007/s00220-026-05705-2}.
\end{thebibliography}
\endgroup

\par\medskip
{\small\noindent\textsc{Beijing Institute of Mathematical Sciences and Applications}\par
\noindent Huairou District, Beijing 101408, China\par}
\end{document}
